\exercisesection{4}

\begin{enumerate}
\def\labelenumi{\arabic{enumi}.}
\tightlist
\item
  \begin{enumerate}
  \def\labelenumii{(\alph{enumii})}
  \tightlist
  \item
    Let G be a totally ordered group and M a major set in \(G_{+}\) not containing 0. Show that there exists a greatest isolated subgroup H of G not meeting M.
  \end{enumerate}
\end{enumerate}

\begin{enumerate}
\def\labelenumi{(\alph{enumi})}
\setcounter{enumi}{1}
\item
  Let A be a valuation ring and v a valuation associated with A. For an ideal \(p \neq 0\) of A to be prime, it is necessary and sufficient that p correspond to a major set M in the totally ordered value group G of v such that, if H is the greatest isolated subgroup of G not meeting M, M is the complement of \(H_{+}\) in \(G_{+}\) .
\item
  With the notation of (b), for an ideal q of A to be p-primary, it is necessary and sufficient that it satisfy one of the following conditions: either \(H = \{0\}\) (and p is then maximal) or q = p.
\item
  For an ideal \(\mathfrak{a}\) of \(\mathbf{A}\) which is distinct from 0 and \(\mathbf{A}\) to satisfy \(\mathfrak{a}^2 = \mathfrak{a}\) , it is necessary and sufficient that \(\mathfrak{a}\) correspond to a major set \(\mathbf{M}\) such that, if \(\mathbf{H}\) is the greatest isolated subgroup of \(\mathbf{G}\) not meeting \(\mathbf{M}\) , the image \(\mathbf{M}'\) of \(\mathbf{M}\) in the totally ordered group \(\mathbf{G}' = \mathbf{G} / \mathbf{H}\) has no least element; \(\mathfrak{a}\) is then a prime ideal.
\end{enumerate}

\begin{enumerate}
\def\labelenumi{\arabic{enumi}.}
\setcounter{enumi}{1}
\tightlist
\item
  Let G be a totally ordered group.
\end{enumerate}

\begin{enumerate}
\def\labelenumi{(\alph{enumi})}
\item
  For every element a \textgreater{} 0 of G, let \(\mathrm{H}(a)\) be the smallest isolated subgroup containing a, \(\mathrm{H}^{-}(a)\) the greatest isolated subgroup not containing a; show that \(\mathrm{H}(a)/\mathrm{H}^{-}(a)\) is isomorphic to a subgroup of R.
\item
  Conversely, if H is an isolated subgroup of G such that there exists a greatest element (with respect to inclusion) \(H'\) in the set of isolated subgroups of G contained in H and \(\neq H\) , then \(\mathrm{H} = \mathrm{H}(a)\) and \(\mathrm{H}' = \mathrm{H}^{-}(a)\) for all \(a \in H_{+} \cap CH_{+}'\) . Such an isolated subgroup H is called principal and \(H'\) is called its predecessor.
\item
  If \(\mathbf{H}_1, \mathbf{H}_2\) are two distinct isolated subgroups of \(\mathbf{G}\) such that \(\mathbf{H}_1 \subset \mathbf{H}_2\) , show that there exist two isolated subgroups \(\mathbf{H}, \mathbf{H}'\) of \(\mathbf{G}\) such that \(\mathbf{H}_1 \subset \mathbf{H}' \subset \mathbf{H} \subset \mathbf{H}_2\) and \(\mathbf{H} / \mathbf{H}'\) is isomorphic to a non-zero subgroup of \(\mathbf{R}\) .
\item
  Let \(\Sigma\) be the set (totally ordered with respect to inclusion) of isolated subgroups of \(G\) and let \(\Theta \subset \Sigma\) be the set of principal isolated subgroups. Show that \(\Sigma\) is isomorphic to the completion (Set Theory, Chapter III, § 1, Exercise 15) of \(\Theta\) .
\end{enumerate}

\begin{enumerate}
\def\labelenumi{\arabic{enumi}.}
\setcounter{enumi}{2}
\tightlist
\item
  \begin{enumerate}
  \def\labelenumii{(\alph{enumii})}
  \tightlist
  \item
    Let \(\Theta\) be a totally ordered set and for each \(s \in \Theta\) let \(\mathbf{A}_s\) be a subgroup of \(\mathbf{R}\) not reduced to 0. Let
  \end{enumerate}
\end{enumerate}

\[
\Gamma (\Theta , (A _ {s}) _ {s \in \Theta}) = G
\]

be the subgroup of the product \(\biguplus_{s\in\Phi}\mathbf{A}_{s}\) consisting of the functions \(f\) whose support is a well-ordered subset of \(\Theta\) . We define on \(\mathbf{G}\) an ordering compatible with its group structure taking the set \(\mathbf{G}_{+}\) of elements \(\geqslant0\) to be the set consisting of the function 0 and the functions \(f\neq0\) such that \(f(\theta)>0\) for the least element \(\theta\) of the support of \(f\) . Show that \(\mathbf{G}\) is a totally ordered group and that \(\Theta\) is canonically isomorphic to the set of principal isolated subgroups of \(\mathbf{G}\) ordered by the relation \(\supset\) ; moreover, if \(\mathrm{H}(a)\) corresponds to \(s\in\Theta\) under this isomorphism, \(\mathrm{H}(a)/\mathrm{H}^{-}(a)\) is isomorphic to \(\mathbf{A}_{s}\) .

\begin{enumerate}
\def\labelenumi{(\alph{enumi})}
\setcounter{enumi}{1}
\tightlist
\item
  Consider the subgroup \(G'\) of the totally ordered group \(Q \times Q\) (with the lexicographical ordering) generated by the elements \((p_{n}^{-1}, np_{n}^{-1})\) , where \((p_{n})\) is the strictly increasing sequence of prime numbers. Show that the only isolated subgroup \(H'\) of \(G'\) distinct from 0 and \(G'\) is the group \(\{0\} \times Z\) ; but \(G'\) is not isomorphic to a subgroup of the product \(H' \times (G'/H')\) ordered by the lexicographical ordering.
\end{enumerate}

\begin{enumerate}
\def\labelenumi{\arabic{enumi}.}
\setcounter{enumi}{3}
\tightlist
\item
  Let \(A\) be a valuation ring which is not a field.
\end{enumerate}

\begin{enumerate}
\def\labelenumi{(\alph{enumi})}
\item
  Show that for A to be a discrete valuation ring, it is necessary and sufficient that every prime ideal of A be principal.
\item
  Show that, if A is the ring of a valuation of height one, the field of fractions of A is a finitely generated A-algebra.
\end{enumerate}

* 5. Let K be a field and M the set of subrings of K which are not fields. Show that a maximal element of M is a valuation ring of height 1 of its field of fractions L and that K is an algebraic extension of L. The following conditions are equivalent:

(α) M is not empty.

(β) M admits a maximal element.

(γ) There exists a valuation of height 1 on K

(δ) K is not an algebraic extension of a finite prime field (cf.~§ 8, no. 1). *

¶ 6. (a) Let S be a hyperstonian compact space with no isolated point (Integration, Chapter V, § 5, Exercise 14). Let \(\mathcal{C}_{0}(S)\) be the vector space of realvalued functions, finite or otherwise, defined on S, continuous on S and such that the sets \(\overline{f}(-\infty)\) and \(\overline{f}(+\infty)\) are nowhere dense in S (Integration, Chapter II, §1, Exercise 13 (g)). Show that, for every measure \(\mu > 0\) on S, there exist functions f in \(\mathcal{C}_{0}(S)\) whose upper integral is infinite (note that for every nowhere dense closed subset N of S there exists a function f in \(\mathcal{C}_{0}(S)\) which is equal to \(+\infty\) on N; show then that attention may be restricted to the case where the measure \(\mu\) is normal and define f suitably as the upper bound of a sequence of functions in \(\mathcal{C}_{0}(S)\) ).

\begin{enumerate}
\def\labelenumi{(\alph{enumi})}
\setcounter{enumi}{1}
\item
  Let \(\mathbf{G}\) be the ordered subgroup of \(\mathcal{C}_0(S)\) consisting of the functions in \(\mathcal{C}_0(S)\) whose values in \(\mathbf{Z}\) or \(\pm \infty\) . Show that \(\mathbf{G}\) is a complete lattice and is not isomorphic (as an ordered group) to any subgroup of a product group \(\mathbf{R}^{\mathbf{l}}\) .
\item
  Let K be a field and \(A_{0} = K[[X_{s}]]_{s \in S}\) the ring of formal power series with coefficients in K in a family of indeterminates ( \(X_{s}\) ) with the space S as indexing set (Algebra, Chapter IV, § 5, Exercise 1). Let b be the set of formal power series \(P \in A_{0}\) with the following property: there exists a nowhere dense set \(N \subset S\) such that every monomial of P whose coefficient is \(\neq 0\) contains an \(X_{s}\) for which \(s \in N\) . Show that b is an ideal of \(A_{0}\) and that the ring \(A = A_{0}/b\) is an integral domain and completely integrally closed (to prove the latter point, apply Chapter V, § 1, no. 4, Proposition 14, arguing as in Chapter V, § 1, Exercise 10; use also the fact that in S every meagre set is nowhere dense).
\item
  Let \(f\) be an element \(\geqslant 0\) of the group \(G\) and let \(N\) be the nowhere dense set of points where \(f(s) = \pm \infty\) . Let \(p_f\) be the element of \(A\) the image of the formal power series \(\prod_{s \notin N} (1 + X_s)^{f(s)}\) ; the mapping \(f \mapsto p_f\) is an isomorphism of the additive monoid \(G_+\) onto a multiplicative submonoid of \(A\) . Deduce that \(A\) cannot be the intersection of a family of valuation rings of height 1 of its field of fractions (show that this would imply that \(G\) is isomorphic to a subgroup of a product \(R^1\) ). (Cf. Exercises 8 and 9.)
\end{enumerate}

¶ 7. A local integral domain A is said to be of dimension 1 if A is not a field and there is no prime ideal of A distinct from (0) and the maximal ideal m of A. For every ideal a of A, let A: a denote the sub-A-module of the field of fractions of A which is the transporter of a in A (which always contains A). In the following A is assumed to be of dimension 1.

\begin{enumerate}
\def\labelenumi{(\alph{enumi})}
\item
  Show that, if A is strongly Laskerian (Chapter IV, § 2, Exercise 28), then A: m ≠ A. (In the opposite case, note that A: m \(^{r}\) = A for every integer r \textgreater{} 0; on the other hand, every ideal ≠ 0 of A contained in m is m-primary and so in particular is every principal ideal Ax ⊆ m (where x ≠ 0); note finally that m \(^{h}\) ≠ m \(^{k}\) for h ≠ k (Chapter IV, § 2, Exercise 29 (d)).)
\item
  Show that for a strongly Laskerian local integral domain A of dimension 1 the following properties are equivalent:
\end{enumerate}

\((\alpha)\) A is completely integrally closed.

(β) The maximal ideal m is principal.

(γ) Every m-primary ideal is of the form \(m^{k}\) .

(δ) The ideal m is invertible (Chapter II, § 5, no. 6), which is equivalent to m. (A: m) = A.

(ε) A is a discrete valuation ring.

(To show that \((\alpha)\) implies \((\delta)\) , use (a) and observe that if \(\mathfrak{m}.\left(\mathbf{A}:\mathfrak{m}\right)^k = \mathfrak{m}\) for all \(k > 0\) , \(\mathbf{A}\) is not completely integrally closed (cf.~Chapter VII, §1, Exercise 4). To see that \((\delta)\) implies \((\gamma)\) , note that for every ideal \(\mathfrak{a} \neq 0\) contained in \(\mathfrak{m}\) we may write \(\mathfrak{a} = \mathfrak{ma}_1\) , where \(\mathfrak{a}_1 \neq 0\) and use Chapter IV, §2, Exercise 29 (d). Finally, to see that \((\gamma)\) implies \((\beta)\) , observe that \(\mathfrak{m} \neq \mathfrak{m}^2\) and that \((\gamma)\) implies that \(\mathfrak{m} / \mathfrak{m}^2\) is a 1-dimensional vector \((\mathbf{A} / \mathfrak{m})\) -space.)

\begin{enumerate}
\def\labelenumi{(\alph{enumi})}
\setcounter{enumi}{2}
\item
  Show, in the notation of Chapter III, § 3, Exercise 15 (b), that the local ring \(\mathbf{A}_{\mathfrak{m}}\) is a Noetherian integral domain of dimension 1 but is not integrally closed.
\item
  Let K be an algebraically closed field and A the subring of the field \(\mathbf{K}(\mathbf{X},\mathbf{Y})\) of rational functions in two indeterminates over K consisting of the elements \(f\in\mathbf{K}(\mathbf{X},\mathbf{Y})\) such that 0 is substitutable for X in f and \(f(0,\mathbf{Y})\) belongs to K. Show that A is a strongly Laskerian integrally closed local integral domain of dimension 1 but that it is not completely integrally closed. (If B is the polynomial ring \(\mathbf{K}[\mathbf{X},\mathbf{Y}]\) , p the prime ideal BX of B, note that A is the subring of \(B_{p}\) equal to \(K+pB_{p}\) .) (*)
\end{enumerate}

¶ 8. Let A be a local integral domain of dimension 1 (Exercise 7) and K its field of fractions. Show that for A to be the intersection of valuation rings (of K) of height 1, it is necessary and sufficient that A be completely integrally closed. Prove successively that:

\begin{enumerate}
\def\labelenumi{(\alph{enumi})}
\item
  A subring of K containing A and the inverse of an element \(\neq0\) of the maximal ideal m of A is equal to K (use the fact that every ideal of A distinct from 0 and A is m-primary).
\item
  If \(\mathbf{B} \supset \mathbf{A}\) is a subring of \(\mathbf{K}\) distinct from \(\mathbf{K}\) , there exists a valuation ring \(\mathbf{V}\) of height 1 such that \(\mathbf{B} \subset \mathbf{V}\) (use (a)).
\item
  Let \(z \in \mathbf{K} - \mathbf{A}\) ; show that \(zm \subset m\) is impossible arguing as in Exercise 7 (b). Deduce that there exists a valuation ring \(V\) of height 1 such that \(A \subset V\) and \(z \notin V\) . (If \(zm \subset A\) , use (a); otherwise, there is an \(x \in m\) such that
\end{enumerate}

\[
x y = z \in \mathrm{K} - \mathrm{A};
\]

observe that \(z \notin A[z^{-1}]\) and use (b) to show the existence of V.)

¶ 9. Let A be a Laskerian completely integrally closed domain (Chapter IV, §2, Exercise 23). Suppose further, that, for all \(x \neq 0\) in A, the prime ideals \(\mathfrak{p}_i\) weakly associated with A/Ax (Chapter IV, § 1, Exercise 17) are all of height 1, that is to say that for all i, \(p_{i}\) contains no prime ideal other than itself and 0 (cf.~Chapter VII, § 1, no. 6). Show that under these conditions A is an intersection of valuation rings of height 1. (Show first that A is the intersection of all the rings \(A_{p}\) , where p runs the set of prime ideals of height 1, using Exercise 17 (i) of Chapter IV, § 1. Then prove that these local rings \(A_{p}\) are completely integrally closed, using condition ( \(LA_{I}\) ) of Chapter IV, § 2, Exercise 23. Finally apply Exercise 8.) The above hypotheses are in particular satisfied when A is a Laskerian completely integrally closed domain in which every ideal admits a single primary decomposition (Chapter IV, § 2, Exercise 26).

¶ 10. Let A be a ring in which the set of principal ideals is totally ordered by inclusion.

\begin{enumerate}
\def\labelenumi{(\alph{enumi})}
\item
  Show that \(\mathbf{A}\) is a local ring, in which the nilradical \(\mathfrak{N}\) is a prime ideal. Show that, either \(\mathfrak{N}^2 = \mathfrak{N}\) , or \(\mathfrak{N}\) is nilpotent; the ring \(\mathbf{A} / \mathfrak{N}\) is a valuation ring.
\item
  If \(\mathcal{L}\) is the totally ordered set of principal ideals of A, show that the set of ideals of A is totally ordered by inclusion and is isomorphic to the completion of the set \(\mathcal{L}\) (Set Theory, Chapter III, § 1, Exercise 15).
\item
  Suppose that \(\mathfrak{N}^2 = 0\) ; then \(\mathfrak{N}\) may be considered as a module over the valuation ring \(V = A / \mathfrak{N}\) ; let \(K\) be the field of fractions of \(V\) and \(\Gamma\) the order group of a valuation \(v\) on \(K\) corresponding to \(V\) ; show that there exists in \(\Gamma\) two major sets \(M \subset M'\) such that \(\mathfrak{N}\) is a \(V\) -module isomorphic to \(\mathfrak{a}(M') / \mathfrak{a}(M)\) (notation of §3, no. 4).
\item
  Conversely, given a valuation ring V with field of fractions K and order group \(\Gamma\) and two major sets \(M \subset M'\) in \(\Gamma\) , let \(Q = a(M') / a(M)\) and define on the product additive group \(A = V \times Q\) a multiplication by setting
\end{enumerate}

\[
(z, t) (z ^ {\prime}, t ^ {\prime}) = (z z ^ {\prime}, z t ^ {\prime} + z ^ {\prime} t);
\]

show that in A the set of principal ideals is totally ordered by inclusion; the set N of ordered pairs \((0, t)\) where \(t \in Q\) is the nilradical of A, \(N^{2} = 0\) and A/N is isomorphic to V.

\begin{enumerate}
\def\labelenumi{(\alph{enumi})}
\setcounter{enumi}{4}
\item
  For every prime number \(p\) , the ring \(A = \mathbf{Z} / p^2\mathbf{Z}\) is such that the set of principal ideals of \(A\) is totally ordered by inclusion and the nilradical \(\mathfrak{N}\) of \(A\) is such that \(\mathfrak{N}^2 = 0\) ; but show that the ring \(A\) is not isomorphic to the ring constructed (starting with \(V = A / \mathfrak{N}\) and \(Q = \mathfrak{N}\) ) by the method in (d) (note that \(A\) contains no subring isomorphic to \(Z / p\mathbf{Z}\) ).
\item
  Show that for every prime number \(p\) , the algebra \(\mathbf{A}\) of the group \(\mathbf{U}_p\) (Algebra, Chapter VII, §2, Exercise 3) with respect to the prime field \(\mathbf{F}_p\) is a ring in which the principal ideals form a set which is totally ordered by inclusion and the nilradical \(\mathfrak{N}\) satisfies \(\mathfrak{N}^2 = \mathfrak{N}\) , but \(\mathbf{A}\) is not isomorphic to the quotient of a valuation ring by an ideal.
\end{enumerate}

\begin{enumerate}
\def\labelenumi{\arabic{enumi}.}
\setcounter{enumi}{10}
\tightlist
\item
  Let K be a field with a valuation v of height 1.
\end{enumerate}

\begin{enumerate}
\def\labelenumi{(\alph{enumi})}
\tightlist
\item
  Let \(\mathbf{P}(\mathbf{X}) = a_0\mathbf{X}^n +a_1\mathbf{X}^{n - 1} + \dots +a_n\) be a polynomial in \(\mathrm{K[X]}\) of degree \(n > 1\) and such that \(a_{n} \neq 0\) ; show that there exists a strictly increasing sequence \((i_{k})_{0 \leqslant k \leqslant r}\) of integers in the interval \([0, n]\) such that: (1) \(i_{0} = 0\) , \(i_{r} = n\) ; (2) \(v(a_{i_{k}})\) is finite for \(0 \leqslant k \leqslant r\) ; (3) for every index \(j\) such that \(0 \leqslant j \leqslant n\) distinct from the \(i_{k}\) , such that \(v(a_{j})\) is finite, the point
\end{enumerate}

\[
(j, v (a _ {j})) \in \mathbf {R} ^ {2}
\]

lies above the line passing through the points \((i_{k}, v(a_{i_{k}}))\) and \((i_{k+1}, v(a_{i_{k+1}}))\) and strictly above that line if \(j < i_{k}\) or \(j > i_{k+1}\) . The union of the segments joining the points \((i_{k}, v(a_{i_{k}}))\) and \((i_{k+1}, v(a_{i_{k+1}}))\) is called the Newton polygon of P, the above segments are called the sides and the points \((i_{k}, v(a_{i_{k}}))\) the vertices of the polygon.

\begin{enumerate}
\def\labelenumi{(\alph{enumi})}
\setcounter{enumi}{1}
\item
  Suppose that all the zeros of P belong to K. Show that for the valuations of all the zeros of P to be the same, it is necessary and sufficient that r = 1 (in other words, that the Newton polygon reduce to a single side). (To show that the condition is sufficient, consider the Newton polygon of a product \(P_{1}P_{2}\) where all the zeros of \(P_{1}\) are invertible in the ring of v, whilst all those of \(P_{2}\) belong to the ideal of v.)
\item
  Suppose that all the zeros of \(\mathbf{P}\) belong to \(\mathbf{K}\) ; form the Newton polygon of \(\mathbf{P}\) and write.
\end{enumerate}

\[
\rho_ {k} = i _ {k + 1} - i _ {k}, \quad \sigma_ {k} = (v (a _ {i _ {k + 1}}) - v (a _ {i _ {k}})) / \rho_ {k}.
\]

Show that, for \(0 \leqslant k \leqslant r - 1\) , P admits exactly \(\rho_k\) zeros (counted with their orders of multiplicity) whose valuations are all equal to \(\sigma_k\) (use (b) and argue by induction on \(r\) ).

\begin{enumerate}
\def\labelenumi{(\alph{enumi})}
\setcounter{enumi}{3}
\tightlist
\item
  Generalize to the case of any valuation v (embed the order group \(\Gamma\) of v in the vector Q-space \(\Gamma_{(Q)}\) , which has a natural totally ordered group structure).
\end{enumerate}

\exercisesection{5}

¶ 1. Let A be an integral domain, K its field of fractions and T a linear topology on A (Chapter III, § 2, Exercise 21).

\begin{enumerate}
\def\labelenumi{(\alph{enumi})}
\item
  For the neighbourhoods of 0 under \(\mathcal{T}\) to constitute a fundamental system of neighbourhoods of 0 for some topology \(\mathcal{T}_{\mathrm{K}}\) compatible with the ring structure on K, it is necessary and sufficient that \(\mathcal{T}\) be the topology \(\mathcal{T}_{u}(\mathrm{A})\) (Chapter III, §2, Exercise 24); then A is a bounded subset with respect to \(\mathcal{T}_{\mathrm{K}}\) and \(\mathcal{T}_{\mathrm{K}}\) is a locally bounded Hausdorff topology (General Topology, Chapter III, §6, Exercises 12 and 20 (e)). For K to be complete (resp. linearly compact. resp. strictly linearly compact (Chapter III, §2, Exercise 21)) with \(\mathcal{T}_{\mathrm{K}}\) , it is necessary and sufficient that A be so with \(\mathcal{T}\) .
\item
  For the topology \(\mathcal{T}_{\mathbf{K}}\) (where \(\mathcal{T} = \mathcal{T}_u(\mathbf{A})\) ) to be compatible with the field structure on \(\mathbf{K}\) , it is necessary and sufficient that the Jacobson radical \(\Re(\mathbf{A})\) of \(\mathbf{A}\) be \(\neq 0\) .
\end{enumerate}

¶ 2. Let K be a (commutative) field and T a non-discrete Hausdorff topology on K compatible with the ring structure on K. For the topology T to be defined by a valuation on K or an absolute value on K, it is necessary and sufficient that T be locally retrobounded (General Topology, Chapter III, § 6, Exercise 22. If there exist in K topologically nilpotent elements, use Exercise 22 (d) of General Topology, Chapter III, § 6 and Exercise 13 of General Topology, Chapter IX, § 3. In the opposite case, use Exercise 22 (f) of General Topology, Chapter III, § 6).

¶ 3. Let A be a Noetherian integral domain, K its field of fractions, \(\mathcal{T}_{u}\) the topology \(\mathcal{T}_{u}(\mathrm{A})\) and \(\mathcal{T}_{\mathrm{K}}\) the corresponding topology on K (Exercise 1).

\begin{enumerate}
\def\labelenumi{(\alph{enumi})}
\item
  If A is a Zariski ring with Jacobson radical \(\mathfrak{r} \neq 0\) and A is complete with the r-adic topology, show that A is complete with the topology \(\mathcal{T}_u\) (General Topology, Chapter III, § 3, no. 5, Corollary 2 to Proposition 9).
\item
  Suppose that \(\mathcal{T}_{\mathbf{K}}\) is not discrete and is defined by a valuation \(v\) on \(\mathbf{K}\) . Show that \(v\) is a discrete valuation and that \(\mathbf{A}\) is an open subring of the ring \(\mathbf{V}\) of \(v\) ; obtain the converse. (Using the hypothesis, which implies \(\mathbf{A} \neq \mathbf{K}\) , and Proposition 1 of §4, no. 1, we may assume that \(\mathbf{A} \subset \mathbf{V}\) . Using the fact that \(\mathbf{A}\) is open under \(\mathcal{T}_{\mathbf{K}}\) , show that \(\mathbf{V}\) is a finitely generated A-module and conclude using §3, no. 5, Proposition 9. For the converse, observe that \(\mathbf{A}\) is a local ring of which the maximal ideal \(m\) and (0) are the only prime ideals and in which every ideal \(a \neq 0\) is therefore \(m\) -primary). The integral closure of \(\mathbf{A}\) is then \(\mathbf{V}\) . Give an example where \(\mathbf{A} \neq \mathbf{V}\) (take \(\mathbf{V}\) to be a ring of formal power series \(k[[T]]\) , where \(k\) is a field).
\item
  Deduce an example of a non-discrete complete topological field whose topology is not locally retrobounded (use (b) and Exercise 2).
\end{enumerate}

\begin{enumerate}
\def\labelenumi{\arabic{enumi}.}
\setcounter{enumi}{3}
\tightlist
\item
  Let \(v\) be a valuation on a field \(K\) , A the ring of \(v\) and \(m\) its maximal ideal.
\end{enumerate}

\begin{enumerate}
\def\labelenumi{(\alph{enumi})}
\item
  For the topology defined by v on A to be identical with the m-adic topology, it is necessary and sufficient that A be a field or a discrete valuation ring.
\item
  For the topology defined by \(v\) to make \(A\) into a strictly linearly compact ring (Chapter III, §2, Exercise 21), it is necessary and sufficient that \(A\) be a field or a discrete valuation ring (use (a) and Chapter III, §2, Exercise 22 (a)).
\end{enumerate}

\begin{enumerate}
\def\labelenumi{\arabic{enumi}.}
\setcounter{enumi}{4}
\tightlist
\item
  \begin{enumerate}
  \def\labelenumii{(\alph{enumii})}
  \tightlist
  \item
    Let K be a field, v a valuation on K and \(\Gamma\) the order group of v. For K to be linearly compact (Chapter III, §2, Exercise 15) with the topology \(T_{v}\) , it is necessary and sufficient that, for every well-ordered subset B of \(\Gamma\) and every family \((a_{\beta})_{\beta\in\mathcal{B}}\) of elements of K such that, for \(\lambda<\mu<v\) ,
  \end{enumerate}
\end{enumerate}

\[
v (a _ {\lambda} - a _ {\mu}) <   v (a _ {\mu} - a _ {\nu}),
\]

there exists an element \(a \in \mathbf{K}\) such that \(v(a - a_{\lambda}) = v(a_{\lambda} - a_{\mu})\) for every ordered pair of indices \(\lambda, \mu\) such that \(\mu > \lambda\) . (Use Exercise 4 of Set Theory,

Chapter III, § 2). If this is so, the ring A of the valuation \(v\) is also linearly compact with the discrete topology.

\begin{enumerate}
\def\labelenumi{(\alph{enumi})}
\setcounter{enumi}{1}
\tightlist
\item
  Show that the field \(S(\Gamma, k) = K\) defined in §3, Exercise 2 is linearly compact with \(T_{v}\) . Take \(\Gamma\) to be the totally ordered group Q of rational numbers and consider the subring \(K_{0}\) of K consisting of the \(x = (x_{\alpha})\) such that the set of \(\alpha \in Q\) for which \(x_{\alpha} \neq 0\) is finite or is the set of points of a strictly increasing sequence tending to \(+\infty\) . Show that \(K_{0}\) is a field which is complete with respect to the valuation \(v_{0}\) induced by the canonical valuation v on K and that \(v_{0}\) has the same value group and residue field as v but that \(K_{0}\) is not linearly compact (cf.~§10, Exercise 2).
\end{enumerate}

¶ 6. (a) Let A be a valuation ring, M an A-module and M' a submodule of M. Show that for M' to be a pure submodule of M (Chapter I, § 2, Exercise 24), it is necessary and sufficient that, for all \(\alpha \in A\) , M' ∩ ( \(\alpha\) M) = \(\alpha\) M'.

\begin{enumerate}
\def\labelenumi{(\alph{enumi})}
\setcounter{enumi}{1}
\item
  Let \(\mathbf{M}'\) be a pure submodule of \(\mathbf{M}\) such that, if \(\mathbf{M}'\) is given the topology for which the \(\alpha\mathbf{M}'\) (where \(\alpha \in \mathbf{A}, \alpha \neq 0\) ) form a fundamental system of neighbourhoods of \(0\) , \(\mathbf{M}'\) is linearly compact. Show that for every element \(x \in \mathbf{M}\) there exists \(y' \in \mathbf{M}'\) such that \(x' = x + y'\) has the following property: for all \(\alpha \in \mathbf{A}\) such that \(x + \mathbf{M}' \in \alpha(\mathbf{M} / \mathbf{M}')\) , \(x' \in \alpha\mathbf{M}\) . (For each of the \(\alpha \in \mathbf{A}\) such that \(x + \mathbf{M}' \in \alpha(\mathbf{M} / \mathbf{M}')\) , consider the subset \(\mathrm{S}_{\alpha}\) of \(\mathbf{M}'\) consisting of the \(y_{\alpha}'\) such that \(x + y_{\alpha}' \in \alpha\mathbf{M}\) .)
\item
  Suppose that A is a linearly compact valuation ring; let K be the field of fractions of A. Show that, if M is a torsion-free A-module such that \(\mathbf{M}_{(\mathbf{K})}\) admits a countable basis, M is the direct sum of a countable family of A-modules of rank 1. (Consider M as the union of an increasing sequence \((\mathbf{M}_{i}^{\prime})\) of pure submodules such that \(M_{i}^{\prime}\) is of rank i and for each i apply (b) to \(M_{i}^{\prime}\) and its submodule \(M_{i-1}^{\prime}\) .
\item
  With the same hypotheses on A and M as in (c), let N be a pure submodule of M of finite rank. Show that N is a direct factor of M (observe that every A-module which is the direct sum of a finite number of A-modules of rank 1 is linearly compact and use (b)).
\item
  With the same hypotheses on M and \(M'\) as in (b), show that for all \(x \in M\) there exists \(y' \in M'\) such that the annihilator of \(x' = x + y'\) is equal to the annihilator of the element \(x + M' \in M/M'\) . (Let a be the annihilator of \(x + M'\) ; for all \(\alpha \in a\) , let \(T_{\alpha}\) be the subset of \(M'\) consisting of the \(y'\) such that \(\alpha(x + y') = 0\) ; show that the intersection of the \(T_{\alpha}\) is non-empty.)
\item
  Suppose that A is a valuation ring such that for every ideal \(a \neq 0\) the A-module A/a is linearly compact (with the discrete topology). Show that every finitely generated torsion A-module M is the direct sum of a finite number of monogenous A-modules. (If \((z_{i})_{1 \leqslant i \leqslant n}\) is a system of generators of M, argue by induction on n, considering one of the \(z_{i}\) whose annihilator is the least and noting that the submodule of M which it generates is pure.)
\end{enumerate}

\begin{enumerate}
\def\labelenumi{\arabic{enumi}.}
\setcounter{enumi}{6}
\tightlist
\item
  Let K be a field and v a discrete valuation on K such that K is complete with \(T_{v}\) : let A and p denote the ring and ideal of v and U = A - p the set of invertible elements of A. Let u be an isomorphism of K onto a subfield of K.
\end{enumerate}

\begin{enumerate}
\def\labelenumi{(\alph{enumi})}
\item
  Show that \(u(\mathfrak{p}) \subset \mathbf{A}\) and \(u(\mathbf{U}) \subset \mathbf{U}\) . (To prove the first point, observe that for all \(z \in \mathfrak{p}\) the equation \(x^n = 1 + z\) admits a solution in \(\mathbf{A}\) for all \(n > 0\) prime to the characteristic of the residue field of \(v\) , using Hensel's Lemma; if \(v(u(z)) < 0\) , deduce that the integer \(v(u(z))\) would be divisible by every integer \(n > 0\) prime to the characteristic of the residue field of \(v\) . Then deduce the second assertion from the first.)
\item
  Deduce from (a) that, either \(u(\mathbf{K}^*) \subset \mathbf{U}\) , or \(u\) is continuous (consider the image under \(u\) of a uniformizer of \(v\) ).
\item
  Give an example of an algebraically closed field \(\Omega\) such that, taking \(\mathbf{K} = \Omega((\mathbf{X}))\) , \(\mathbf{K}\) is isomorphic to a subfield of \(\mathbf{K}\) contained in \(\mathbf{U} \cup \{0\}\) (cf.~Algebra, Chapter V, § 5, Exercise 13).
\end{enumerate}

¶ 8. (a) Let K be a field, v a valuation on K of height 1 and A its ring. Let H be a compact subset of K (with the topology \(T_{v}\) ) and \(a \neq 0\) a point of K. Show that there exists a polynomial \(f \in K[X]\) without constant term and such that \(f(a) = 1\) , \(f(H) \subset A\) . (Prove that f may be taken to be a polynomial of the form

\[
1 - (1 - a ^ {- 1} \mathrm{X}) (1 - c _ {1} ^ {- 1} \mathrm{X}) ^ {n (1)} \dots (1 - c _ {r} ^ {- 1} \mathrm{X}) ^ {n (r)}
\]

where the \(c_{i}\) are suitably chosen elements of \(\mathbf{H}\) such that \(v(c_{i}) < v(a)\) and the \(n(i)\) are sufficiently large integers \(>0\) .

\begin{enumerate}
\def\labelenumi{(\alph{enumi})}
\setcounter{enumi}{1}
\tightlist
\item
  Let X be a totally disconnected compact space; let the ring \(\mathcal{C}(\mathbf{X};\mathbf{K})\) of continuous mappings from X to K be given the uniform convergence topology. Let B be a subring of \(\mathcal{C}(\mathbf{X};\mathbf{K})\) containing the constants and separating the points of X; show that B is dense in \(\mathcal{C}(\mathbf{X};\mathbf{K})\) (Use (a) and General Topology, Chapter X, §4, Exercise 21 (b).)
\end{enumerate}

¶ 9. Let A be a discrete valuation ring, π a uniformizer of A and K the field of fractions of A; suppose that A is complete with the π-adic topology. The injective A-modules (Algebra, Chapter II, §2, Exercise 11) are identical with the divisible A-modules (Algebra, Chapter VII, §2, Exercise 3) and are direct sums of A-modules isomorphic either to K or to K/A (Algebra, Chapter VII, §2, Exercise 3). Moreover, every monogenous A-module is isomorphic to a submodule of K/A. The A-module Hom \(_{A}\) (M, K/A) is called the (algebraic) toric dual of an A-module M and denoted by M*; the canonical mapping c \(_{M}\) : M → M** is injective (Algebra, Chapter II, §2, Exercise 13 (b)). For every submodule N of M the submodule N \(^{0}\) of M* consisting of the u such that u(x) = 0 is called the orthogonal of N in M*; the dual of M/N is canonically identified with N \(^{0}\) and the dual of N with M*/N \(^{0}\) . The toric dual of a direct limit lim M \(_{\alpha}\) is canonically isomorphic to the inverse limit lim M \(_{\alpha}^{*}\) .

\begin{enumerate}
\def\labelenumi{(\alph{enumi})}
\tightlist
\item
  We know that the A-modules of rank one (Algebra, Chapter VII, § 4,
\end{enumerate}

Exercise 22) are isomorphic to a module of one of the form A, K, K/A or A/ \(\pi^{h}\) A. Show that the algebraic toric duals of K and A/ \(\pi^{h}\) A are respectively isomorphic to K and A/ \(\pi^{h}\) A, that the toric dual of A is isomorphic to K/A and that that of K/A is isomorphic to A (use the knowledge about the sub-A-modules of K and the fact that A is complete (with respect to the dual of K/A)). Deduce that, for every A-module M of finite rank, M* is a module of the same rank and that the canonical homomorphism \(c_{M}\) is bijective.

\begin{enumerate}
\def\labelenumi{(\alph{enumi})}
\setcounter{enumi}{1}
\item
  Let M be an A-module and N a submodule of M of finite rank. Show that the orthogonal of \(N^{0}\) in \(M^{**}\) is identified (by \(c_{M}\) ) with N (use (a)).
\item
  Show that an A-module M which is Noetherian (resp. Artinian) is of finite rank (embed M in its injective envelope (Algebra, Chapter II, §2, Exercise 18)); then M* is Artinian (resp. Noetherian).
\item
  Let \(\sigma(M^{*}, M)\) denote the topology (on \(M^{*}\) ) of pointwise convergence on \(M\) ( \(K/A\) being given the discrete topology). Show that, if \(N\) is a submodule of \(M\) , \(\sigma(N^{0}, M/N)\) is induced on \(N^{0}\) by \(\sigma(M^{*}, M)\) and that \(\sigma(M^{*}/N^{0}, N)\) is the quotient by \(N^{0}\) of the topology \(\sigma(M^{*}, M)\) . (For the second point, note that, if \(P\) is a submodule of \(M\) such that \(N \subset P\) , the dual of \(P/N\) is identified with \(N^{0}/P^{0}\) ). If \(M = \lim M_{\alpha}\) , the topology \(\sigma(M^{*}, M)\) is the inverse limit of the topologies \(\sigma(M_{\alpha}^{*}, \overrightarrow{M_{\alpha}})\) .
\item
  The topologies \(\sigma(\mathbf{K},\mathbf{K})\) and \(\sigma(\mathbf{A},\mathbf{K}/\mathbf{A})\) are the \(\pi\) -adic topologies; the topologies \(\sigma(\mathbf{A}/\pi^{h}\mathbf{A},\mathbf{A}/\pi^{h}\mathbf{A})\) and \(\sigma(\mathbf{K}/\mathbf{A},\mathbf{A})\) are the discrete topologies. Deduce that for every A-module M, the module \(\mathbf{M}^{*}\) with \(\sigma(\mathbf{M}^{*},\mathbf{M})\) is linearly compact (Chapter III, §2, Exercise 15; consider M as a quotient module of a free A-module).
\item
  Let M, N be two A-modules; for every homomorphism \(u: \mathbf{M} \to \mathbf{N}\) , let \(^{t}u\) denote the homomorphism \(\operatorname{Hom}(u, \mathbf{l}_{\mathbf{K/A}})\) from \(\mathbf{N}^{*}\) to \(\mathbf{M}^{*}\) such that
\end{enumerate}

\[
\left(^ {t} u (w)\right) (x) = w (u (x))
\]

for all \(x \in \mathbf{M}\) and all \(w \in \mathbf{N}^*\) ; show that \({}^t u\) is continuous for the topologies \(\sigma(\mathbf{N}^*, \mathbf{N})\) and \(\sigma(\mathbf{M}^*, \mathbf{M})\) . If \(u\) is the endomorphism \(x \mapsto \pi x\) of \(\mathbf{M}\) , \({}^t u\) is the endomorphism \(w \mapsto \pi w\) of \(\mathbf{M}^*\) . For every submodule \(\mathbf{P}\) of \(\mathbf{M}\) , \((u(\mathbf{P}))^0 = {}^t u^{-1}(\mathbf{P}^0)\) .

¶ 10. The hypotheses and notation are the same as in Exercise 9.

\begin{enumerate}
\def\labelenumi{(\alph{enumi})}
\item
  If M is a discrete topological A-module, show that M is a torsion module. If further M is linearly compact, show that M is Artinian (if N is the kernel of the endomorphism \(x \mapsto \pi x\) , observe that N is linearly compact and discrete and can be considered as a vector (A/ \(\pi\) A)-space; then use Exercise 20 (d) of Chapter II, § 2).
\item
  Deduce from (a) that every linearly compact topological A-module is strictly linearly compact (Chapter II, § 2, Exercise 19).
\item
  The topological toric dual of a topological A-module M is the submodule \(M'\) of the algebraic toric dual \(M^{*}\) consisting of the continuous homomorphisms from M to K/A (the latter being given the discrete topology). If the topology on M is linear (Chapter II, § 2, Exercise 14) and N is a closed submodule of M, show that the topological toric dual of M/N is identified with \(N^{0} \cap M'\) and that of N with \(M' / (N^{0} \cap M')\) (to determine the dual of N, note that, for every continuous homomorphism u from N to K/A, there exists an open submodule U of M such that \(u(x) = 0\) in \(N \cap U\) , then use Exercise 9).
\item
  For a topological A-module M of finite rank, the topological toric dual \(M'\) is equal to the algebraic toric dual \(M^{*}\) (reduce it the case of modules of rank one).
\item
  Let M be a Hausdorff topological A-module whose topology is linear; show that for all \(x \neq 0\) in M, there exists \(u \in M'\) such that \(u(x) \neq 0\) (note that there exists an open submodule U of M such that \(x \notin U\) ); in other words, the canonical mapping \(\mathbf{M} \to (\mathbf{M}')^*\) is injective. Deduce that, if N is a closed submodule of M, \(N^0 \cap M'\) is dense in \(N^0\) with the topology \(\sigma(\mathbf{M}^*, \mathbf{M})\) (use (d)), and consequently \(N = M \cap (N^0 \cap M')^0\) in \(M^{**}\) .
\end{enumerate}

Show that M is dense in \((\mathbf{M}^{\prime})^{*}\) with the topology \(\sigma((\mathbf{M}^{\prime})^{*}, \mathbf{M}^{\prime})\) .

\begin{enumerate}
\def\labelenumi{(\alph{enumi})}
\setcounter{enumi}{5}
\item
  Let M be a linearly compact topological A-module; show that the canonical injection \(M \rightarrow (M')^{*}\) is a bijection and that the topology on M (identified with \((M')^{*}\) ) is equal to \(\sigma(M, M')\) . (Observe that, if U is an open submodule of M, M/U is Artinian by (a) and therefore \(U^{0}\) is Noetherian (Exercise 9 (c)) and deduce the identities of the topologies considered on M; complete with the aid of (e).)
\item
  Let M, N be two topological A-modules whose topologies are linear; for every continuous homomorphism \(u \colon \mathbf{M} \to \mathbf{N}\) , \(^t u(\mathbf{N}') \subset \mathbf{M}'\) and \(^t u\) is also used (by an abuse of notation) to denote the linear mapping from \(\mathbf{N}'\) to \(\mathbf{M}'\) with the same graph as \(^t u\) . If M and N are Hausdorff, show that the restriction to M of \(^t (^t u)\) coincides with \(u\) (M and N being respectively considered as canonically embedded in \((\mathbf{M}')^*\) and \((\mathbf{N}')\) ); moreover, for every submodule Q of N,
\end{enumerate}

\[
(\bar {u} ^ {1} (\mathbf {Q})) ^ {0} \cap \mathbf {M} ^ {\prime} = ^ {t} u (\mathbf {Q} ^ {0})
\]

(use (e)).

\begin{enumerate}
\def\labelenumi{(\alph{enumi})}
\setcounter{enumi}{7}
\tightlist
\item
  Let M be a linearly compact A-module; deduce from (g) and Exercise 9 (f) that for M to be torsion-free, it is necessary and sufficient that M' be divisible and that for M to be divisible, it is necessary and sufficient that M' be torsion-free.
\end{enumerate}

\exercisesection{6}

¶ 1. Every element z of the p-adic field \(Q_{p}\) (p a prime number) may be written uniquely as \(p^{h}\sum_{k=0}^{\infty}c_{k}p^{k}\) , where \(h\in Z\) , \(c_{0}\neq0\) , \(0\leqslant c_{k}<p\) for all \(k\geqslant0\) (``p-adic expansion'' of z).

\begin{enumerate}
\def\labelenumi{(\alph{enumi})}
\tightlist
\item
  Show that, in order that \(z \in \mathbf{Q}\) , it is necessary and sufficient that there exist two integers \(m \geqslant 0, n \geqslant 1\) such that \(c_{k+n} = c_k\) for all \(k \geqslant m\) . (Observe that, if \(z = a / b \in \mathbf{Q}\) where \(b\) is not divisible by \(p\) , then
\end{enumerate}

\[
a - b \sum_ {k = 0} ^ {n} c _ {k} p ^ {k} = a _ {n + 1} p ^ {n + 1}
\]

where \(a_{n+1}\) is an integer and the sequence \(|a_k|\) (ordinary absolute value) is bounded.)

\begin{enumerate}
\def\labelenumi{(\alph{enumi})}
\setcounter{enumi}{1}
\tightlist
\item
  Suppose that the increasing sequence \((k_n)\) of integers \(k\) such that \(c_k \neq 0\) is such that \(\limsup_{n \to \infty} (k_{n+1}/k_n) = +\infty\) . Show that \(z\) is transcendental over \(\mathbf{Q}\) . (For \(x \in \mathbf{Q}\) , let \(|x|\) and \(|x|_p\) respectively denote the usual absolute value and a \(p\) -adic absolute value. Suppose that \(\mathrm{P} \in \mathbf{Z}[\mathrm{X}]\) is an irreducible polynomial such that \(\mathrm{P}(z) = 0\) ; if \(z_n = p^n \sum_{k=0}^{\infty} c_k p^k\) , show that \(|\mathrm{P}(z_n)/(z - z_n)|_p\) tends to a limit \(\neq 0\) as \(n\) tends to \(+\infty\) , using Taylor's formula; then obtain a contradiction from the hypothesis considering the usual absolute values \(|\mathrm{P}(z_n)|\) .)
\end{enumerate}

\begin{enumerate}
\def\labelenumi{\arabic{enumi}.}
\setcounter{enumi}{1}
\tightlist
\item
  Let D be a non-commutative field of characteristic \(\neq2\) and Z its centre. Suppose that every commutative subfield K of D containing Z is of degree \(\leqslant2\) over Z.
\end{enumerate}

\begin{enumerate}
\def\labelenumi{(\alph{enumi})}
\item
  Show without using Lemma 1 that \([\mathbf{D}:\mathbf{Z}] = 4\) . (Let \(a\in \mathbf{D} - \mathbf{Z}\) be such that \(a^2\in \mathbf{Z}\) and let \(\sigma (x) = axa^{-1}\) for all \(x\in \mathbf{D}\) ; note that \(\mathbf{D}\) decomposes as a direct sum of two vector subspaces over \(\mathbf{Z}\) , \(\mathrm{D}_{+}\) and \(\mathrm{D}_{-}\) such that \(\sigma (x) = x\) on \(\mathrm{D}_{+}\) , \(\sigma (x) = -x\) on \(\mathrm{D}_{-}\) ; note also that \(\mathrm{D}_{+}\) is a subfield of \(\mathbf{D}\) and \(\mathrm{D}_{-}\) is a 1-dimensional vector space over \(\mathrm{D}_{+}\) ; finally, show that \(\mathrm{D}_{+}\) cannot be distinct from \(\mathbf{Z}(a)\) .)
\item
  Show that \(\mathbf{D}\) is a quaternion algebra over \(\mathbf{Z}\) . (Form a basis of \(\mathbf{D}\) over \(\mathbf{Z}\) with the aid of (a).)
\end{enumerate}

\exercisesection{7}

\begin{enumerate}
\def\labelenumi{\arabic{enumi}.}
\item
  Let \((\phi_{t})_{t\in I}\) be a family of places of a field \(K\) taking their values in a finite number of fields; suppose further that for some \(x\in K\) , the set of \(\phi_t(x)\) is finite. Show that there exists a polynomial \(f(X)\) of the form (1) of no. 1 such that \(f(x)\neq 0\) and that the element \(z = f(x)^{-1}\) enjoys the following properties: (i) \(\phi_t(x) = \infty\) implies \(\phi_t(xz) = 0\) and \(\phi_t(z) = 0\) ; (ii) \(\phi_t(x)\neq \infty\) implies \(\phi_t(xz)\neq \infty\) and \(\phi_t(z)\neq 0\) . (Same proof as for Lemma 1).
\item
  With the hypotheses and notation of Proposition 1 of no. 1, let q be a prime ideal of B; show that \(B_{q}\) is a valuation ring.
\item
  Let \(\mathbf{A}_i\) ( \(1 \leqslant i \leqslant n\) ) be pairwise independent valuation rings of a field \(\mathbf{K}\) and let \(\mathbf{A} = \bigcap_{i} \mathbf{A}_i\) . For all \(i\) , let \(a_i\) be an ideal \(\neq 0\) of \(\mathbf{A}_i\) and write \(a = \bigcap_{i} a_i\) ;
\end{enumerate}

show that, for all \(i\) , \(\mathfrak{a}_i = \mathbf{A}_i\mathfrak{a}\) . Conversely, if \(\mathfrak{b}\) is an ideal \(\neq 0\) of \(\mathbf{A}\) and \(\mathfrak{b}_i = \mathbf{A}_i\mathfrak{b}\) , then \(\mathfrak{b} = \bigcap_i \mathfrak{b}_i\) (use Theorem 1 of no. 2).

\exercisesection{8}

\begin{enumerate}
\def\labelenumi{\arabic{enumi}.}
\item
  Let K be a field and A = K{[}{[}X, Y, Z{]}{]} the ring of formal power series in three indeterminates over K; let \(v'\) (resp. \(v''\) ) be the valuation on A with values in the group \(Z \times Z\) ordered lexicographically, such that \(v'(X) = (1, 0)\) , \(v'(Y) = (0, 1)\) , \(v'\) is improper in K{[}{[}Z{]}{]} (resp. \(v''(X) = (1, 0)\) , \(v''\) is improper in K{[}{[}Y{]}{]}, \(v''(Z) = (0, 1)\) ). Let \(\sigma\) be the automorphism of A leaving invariant K and X and such that \(\sigma(Y) = Z\) , \(\sigma(Z) = Y\) ; if B is the subring of A consisting of the elements invariant under \(\sigma\) and E (resp. F) the field of fractions of A (resp. B), the valuations \(v'\) and \(v''\) (extended canonically to E) have the same restriction v to F, F is complete with the topology defined by v and E is a quadratic extension of F; the two valuations \(v'\) , \(v''\) on E are dependent but not equivalent.
\item
  Let \(K_{0}\) be the field obtained by adjoining to the 2-adic field \(Q_{2}\) the roots of all the polynomials \(X^{2^{n}} - 2\) ; let v be the unique valuation on \(K_{0}\) extending the 2-adic valuation and let K be the completion of \(K_{0}\) with respect to v. Show that the polynomial \(X^{2} - 3\) is irreducible in K{[}X{]}; let \(K'\) be the field of roots of this polynomial and let \(v'\) be the extension of v to \(K'\) ; then
\end{enumerate}

\[
n = \left[ \mathrm{K} ^ {\prime}: \mathrm{K} \right] = 2, \quad e (v ^ {\prime} / v) = f (v ^ {\prime} / v) = 1.
\]

\begin{enumerate}
\def\labelenumi{\arabic{enumi}.}
\setcounter{enumi}{2}
\tightlist
\item
  Let k be the field of rational functions \(\mathbf{F}_{p}(\mathbf{X}_{n})_{n\in\mathbf{N}}\) in an infinity of indeterminates over the prime field \(F_{p}\) and let \(K = k(\mathbf{U}, \mathbf{V})\) be the field of rational functions in two indeterminates over k.
\end{enumerate}

\begin{enumerate}
\def\labelenumi{(\alph{enumi})}
\item
  Show that the element \(\mathbf{P}(\mathbf{U}) = \sum_{n=0}^{\infty} \mathbf{X}_n^p \mathbf{U}^{np}\) of the field of formal power series \(k((\mathbf{U}))\) is not algebraic over the field \(k(\mathbf{U})\) of rational functions. The mapping \(\mathbf{F}(\mathbf{U}, \mathbf{V}) \mapsto \mathbf{F}(\mathbf{U}, \mathbf{P}(\mathbf{U}))\) of \(k[\mathbf{U}, \mathbf{V}]\) to \(k((\mathbf{U}))\) can be extended to an isomorphism of \(\mathbf{K}\) onto a subfield of \(k((\mathbf{U}))\) ; the restriction to this subfield of the valuation on \(k((\mathbf{U}))\) equal to the order of the formal power series (§ 3, no. 3, Example 3) is a discrete valuation \(v\) on \(\mathbf{K}\) , whose residue field is \(k\) .
\item
  Let \(\mathbf{K}'\) be the algebraic extension \(\mathbf{K}(\mathbf{V}^{1/p})\) of \(\mathbf{K}\) so that \([\mathbf{K}':\mathbf{K}] = p\) ; if \(v'\) is the unique valuation on \(\mathbf{K}'\) extending \(v\) , show that \(e(v'/v) = f(v'/v) = 1\) . The ring of the valuation \(v'\) is therefore not a finitely generated module over the ring of the valuation \(v\) .
\end{enumerate}

\begin{enumerate}
\def\labelenumi{\arabic{enumi}.}
\setcounter{enumi}{3}
\item
  Let \(k\) be a field, \(\mathbf{K} = k(\mathbf{X}, \mathbf{Y})\) the field of rational functions in two indeterminates over \(k\) and \(v\) the valuation on \(\mathbf{K}\) with values in the group \(\mathbf{Z} \times \mathbf{Z}\) ordered lexicographically, such that \(v(\mathbf{X}) = (0, 1)\) , \(v(\mathbf{Y}) = (1, 0)\) . Let \(\mathbf{K}'\) be the field \(\mathbf{K}(\sqrt{\mathbf{X}})\) ; show that \(v\) has a single extension \(v'\) to \(\mathbf{K}'\) and that \(e(v'/v) = 2\) \(f(v^{\prime}/v)=1\) , but the ring of the valuation \(v^{\prime}\) is not a finitely generated module over the ring of the valuation \(v\) .
\item
  Let K be a field, A a valuation ring of K and L a finite algebraic extension of K. Let \(A' \supset A\) be another valuation ring of K; let \(A_i' (1 \leqslant i \leqslant m)\) be the valuation rings of L such that \(A_i' \cap K = A'\) and let \(k_i'\) be their respective residue fields. Let k be the residue field of \(A'\) and \(A''\) the valuation ring of K the canonical image of A; finally let \(A_{ij}' (1 \leqslant j \leqslant n_i)\) be the valuation rings of \(k_i'\) such that \(A_{ij}'' \cap k = A''\) and \(A_{ij}\) the inverse images in \(A_i'\) of the \(A_{ij}'' (1 \leqslant i \leqslant m, 1 \leqslant j \leqslant n_i)\) .
\end{enumerate}

\begin{enumerate}
\def\labelenumi{(\alph{enumi})}
\item
  Show that the \(\mathbf{A}_{ij}\) are valuation rings of \(\mathbf{L}\) which are distinct and such that \(\mathbf{A}_{ij} \cap \mathbf{K} = \mathbf{A}\) and every valuation ring \(\mathbf{B}\) of \(\mathbf{L}\) such that \(\mathbf{B} \cap \mathbf{K} = \mathbf{A}\) is equal to one of the \(\mathbf{A}_{ij}\) .
\item
  Show that
\end{enumerate}

\[
e \left(\mathrm{A} _ {i j} / \mathrm{A}\right) = e \left(\mathrm{A} _ {i} ^ {\prime} / \mathrm{A} ^ {\prime}\right) e \left(\mathrm{A} _ {i j} ^ {\prime \prime} / \mathrm{A} ^ {\prime \prime}\right) \quad \text { and } \quad f \left(\mathrm{A} _ {i j} / \mathrm{A}\right) = f \left(\mathrm{A} _ {i j} ^ {\prime \prime} / \mathrm{A} ^ {\prime \prime}\right).
\]

¶ 6. Let K be a field, v a valuation on K and A the ring of v.

\begin{enumerate}
\def\labelenumi{(\alph{enumi})}
\item
  Suppose that A is Henselian (Chapter III, § 4, Exercise 3) in which case we also say, by an abuse of language, that K is Henselian for v. Then, for every algebraic extension L of K and every valuation \(v'\) on L extending v, the ring \(A'\) of \(v'\) is Henselian (Chapter III, § 4, Exercise 4).
\item
  If K is complete with v and v is of height 1, A is Henselian. Give an example where K is complete with v and v is of height 2, but A is not Henselian (cf.~Exercise 1).
\item
  If K is linearly compact with v, show that K is Henselian. (In the notation condition (H) of Chapter III, § 4, Exercise 3, let L denote the set of prime ideals of A (totally) ordered by inclusion and \(\mathcal{L}\) the set of ordered pairs (B, φ), where B is a well ordered subset (with respect to the relation ⊃) of L and φ: p ↦ (Q\_p, Q'\_p) a mapping from B to the set A{[}X{]} × A{[}X{]}, with the following properties: (1) Q\_p and Q'\_p are monic of respective degrees deg( \(\overline{Q}\) ) and deg( \(\overline{Q}'\) ); (2) if p ⊃ q are two elements of B, the coefficients of the polynomials Q\_p - Q\_q and Q'\_p - Q'\_q belong to p; (3) the coefficients of P - Q\_pQ'\_p belong to p; (4) \(\bar{f}(Q_{p}) = \overline{Q}, \bar{f}(Q'_{p}) = \overline{Q}'\) . Define on L an ordering by setting (B, φ) ≤ (B', φ') when B ⊂ B' and φ' is an extension of φ to B'. Prove that L admits a maximal element (B₀, φ₀) and that B₀ has a last element equal to (0): obtain a contradiction from the opposite by considering two cases, according to whether or not B₀ has a last element; in the former case, if p is this last element, consider an element c ∈ A such that v(c) is the least value taken by v on the set of coefficients of P - Q\_pQ'\_p and the least prime ideal p' of A containing c; if p' = p, use Exercise 2 (a) of § 4 and the fact that a quotient of a linearly compact module by a closed submodule is linearly compact and hence complete. If on the contrary B₀ has no last element, use directly the hypothesis that A is linearly compact.)
\item
  In the notation of Exercise 5, show that for A to be Henselian, it is necessary and sufficient that \(\mathbf{A}'\) and \(\mathbf{A}''\) be so. (Note that, if \(\mathbf{P} \in \mathbf{A}'[\mathbf{X}]\) , there exists \(a \in \mathbf{A}\) such that \(a^n\mathbf{P}(\mathbf{X}/a) \in \mathbf{A}[\mathbf{X}]\) (where \(n\) is the degree of \(\mathbf{P}\) ); to verify axiom (H) of Chapter III, § 4, Exercise 3 for \(\mathbf{A}'\) , attention may be confined to the polynomials of \(\mathbf{A}[\mathbf{X}]\) ; use the fact that \(\mathbf{A}' = \mathbf{A}_{\mathfrak{p}}\) , where \(\mathfrak{p}\) is a prime ideal of \(\mathbf{A}\) contained in the maximal ideal \(\mathfrak{m}\) , and note that, if two polynomials in \((\mathbf{A}/\mathfrak{p})[\mathbf{X}]\) are strongly relatively prime, so are their images in \((\mathbf{A}/\mathfrak{m})[\mathbf{X}]\) .)
\item
  Suppose that A is Henselian; then, for every algebraic extension L of K, there exists (up to equivalence) only one valuation on L extending v. (If \(P \in A[X]\) is a monic irreducible polynomial, \(x_{i}\) ( \(1 \leqslant i \leqslant m\) ) the distinct roots of P in its splitting field N, show that, for every valuation \(v'\) on N extending v, the \(v'(x_{i})\) are all equal for \(1 \leqslant i \leqslant m\) .) Conversely, if the valuation ring A has this property, it is Henselian (observe that, in the integral closure \(A'\) of A in L, there can only exist one maximal ideal lying above m (Chapter V, § 1, Exercise 13) and conclude that, if \(P \in A[X]\) is a monic irreducible polynomial its image in \((\mathbf{A}/\mathfrak{m})[\mathbf{X}]\) cannot be a product of two relatively prime polynomials).
\end{enumerate}

\begin{enumerate}
\def\labelenumi{\arabic{enumi}.}
\setcounter{enumi}{6}
\tightlist
\item
  Let K be a field, v a valuation on K, A the ring of v and m the ideal of v. Suppose that A is Henselian; let L be a finite extension of K of degree n; let \(v'\) be the unique valuation on L extending v (Exercise 6 (e)), A' its ring and \(m'\) its ideal. Let x be any element of A'; show that the degree over A/m of the class \(\bar{x}\) of x in A'/m' is a divisor of n and that the order of the class of \(v'(x)\) in \(\Gamma_{v'}/\Gamma_v\) is a divisor of n (consider the minimal polynomial of x over K).
\end{enumerate}

¶ 8. Let K' be a field, B the ring of a valuation v on K', G a finite group of automorphisms of K', K the subfield of K' consisting of the elements invariant under G and A = K ∩ B, which is a valuation ring of K, corresponding to the valuation w = v\textbar K.

\begin{enumerate}
\def\labelenumi{(\alph{enumi})}
\item
  The valuations on \(\mathbf{K}'\) which extend \(w\) are the \(v \circ \sigma\) , where \(\sigma \in \mathcal{G}\) (no. 6, Corollary 1 to Proposition 6) and the integral closure \(\mathbf{A}'\) of \(\mathbf{A}\) in \(\mathbf{K}'\) is the intersection of the \(\sigma \cdot \mathbf{B}\) , where \(\sigma\) runs through \(\mathcal{G}\) . If \(\mathfrak{p}(\mathbf{B})\) is the intersection of \(\mathbf{A}'\) and the maximal ideal \(\mathfrak{m}(\mathbf{B})\) , then \(\sigma \cdot \mathfrak{p}(\mathbf{B}) = \mathfrak{p}(\sigma \cdot \mathbf{B})\) . Show that the decomposition group \(\mathcal{G}^{\mathbb{Z}}(\mathfrak{p}(\mathbf{B}))\) (Chapter V, § 2, no. 2) is the subgroup of \(\mathcal{G}\) consisting of the \(\sigma\) such that \(\sigma \cdot \mathbf{B} = \mathbf{B}\) ; write \(\mathcal{G}^{\mathbb{Z}} = \mathcal{G}^{\mathbb{Z}}(\mathfrak{p}(\mathbf{B}))\) and denote by \(\mathbf{K}^{\mathbb{Z}}\) the subfield of \(\mathbf{K}'\) consisting of the elements invariant under \(\mathcal{G}^{\mathbb{Z}}\) and by \(\mathbf{B}^{\mathbb{Z}}\) the ring of the valuation induced on \(\mathbf{K}^{\mathbb{Z}}\) by \(v\) , which is equal to \(\mathbf{B} \cap \mathbf{K}^{\mathbb{Z}}\) ; recall that the residue fields of \(\mathbf{B}^{\mathbb{Z}}\) and \(\mathbf{A}\) are the same and that the maximal ideal \(\mathfrak{m}(\mathbf{B}^{\mathbb{Z}})\) is generated by \(\mathfrak{m}(\mathbf{A})\mathbf{B}^{\mathbb{Z}}\) (Chapter V, § 2, no. 2, Proposition 4).
\item
  Let \(v^{\mathbf{Z}}\) denote the restriction of \(v\) to \(\mathbf{K}^{\mathbf{Z}}\) and \(\Gamma\) and \(\Gamma^{\mathbf{Z}}\) the order groups of \(w\) and \(v^{\mathbf{Z}}\) . Show that \(\Gamma = \Gamma^{\mathbf{Z}}\) . (Let \(\mathcal{G}^{\mathrm{D}} \supset \mathcal{G}^{\mathbf{Z}}\) be the subgroup of \(\mathcal{G}\) consisting of the \(\sigma \in \mathcal{G}\) such that \(\sigma\) . B is dependent on B (§ 7, no. 2). Argue by induction on the order of \(\mathcal{G}\) . If \((\mathcal{G}: \mathcal{G}^{\mathrm{D}}) > 1\) , let \(\mathbf{K}''\) be the field of invariants of \(\mathcal{G}^{\mathrm{D}}\) and \(\mathbf{A}'' = \mathbf{K}'' \cap \mathbf{B}\) ; show that the order group of \(v|\mathbf{K}''\) is equal to \(\Gamma\) : for this, use the Approximation Theorem (§ 7, no. 2, Theorem 1) in order to prove that, for all \(x \in K''\) , there exists \(y \in K''\) such that \(v(y) = v(x)\) and \(v(y_i) = 0\) for the conjugates \(y_i\) of \(y\) with respect to \(K\) , distinct from \(y\) ; then \(G\) may be replaced by \(G^D\) and the induction hypothesis may be applied. If \(G^D = G\) , let \(B_0\) be the valuation ring of \(K'\) generated by the \(\sigma.B\) , where \(\sigma \in G\) : let \(\bar{K}' = \kappa(B_0)\) be its residue field, \(\pi: B_0 \to \bar{K}'\) the canonical homomorphism, \(\bar{B} = \pi(B)\) a valuation ring of \(\bar{K}'\) and \(\bar{v}\) the corresponding valuation on \(\bar{K}'\) such that \(v = \bar{v} \circ \pi\) in \(B_0\) , so that the order group \(\bar{\Delta}\) of \(\bar{v}\) is a subgroup of the order group \(\Delta\) of \(v\) ; if \(\Delta_0 = \Delta/\bar{\Delta}\) and \(\bar{\omega}: \Delta \to \Delta_0\) is the canonical homomorphism, \(v_0 = \bar{\omega} \circ v\) is a valuation on \(K'\) corresponding to \(B_0\) and invariant under \(G\) . By taking quotients, \(G\) defines a group of automorphisms \(\bar{G}\) of \(\bar{K}'\) ; let \(N\) be the kernel of the canonical homomorphism \(G \to \bar{G}\) ; show that \(N \subset G^Z\) and deduce that, if \(N \neq \{e\}\) , \(G\) may be replaced by \(G/N\) and the induction hypothesis may be applied. Suppose finally that \(N = \{e\}\) ; let \(A_0 = K \cap B_0\) and \(w_0\) be the restriction of \(v_0\) to \(K\) ; show that the order group of \(w_0\) is \(\Delta_0\) and that \(\pi(A_0) = \bar{K}\) is the field of invariants of \(\bar{G}\) (cf.~no. 1, Lemma 2); if \(\bar{w}\) is the restriction to \(\bar{K}\) of \(\bar{v}\) and \(\bar{\Gamma}\) its order group, \(\Delta_0\) is canonically isomorphic to \(\Gamma/\bar{\Gamma}\) . On the other hand, the group \(G^Z(p(\bar{B}))\) is the canonical image of \(G^Z(p(B))\) in \(G^Z\) ; observe that \(G^D \neq G^Z\) and that the induction hypothesis may therefore be applied to prove that \(\bar{\Gamma}\) is equal to the order group \(F^Z(θ)\) of the restriction of \(\bar{v}\) to \(K^Z\). )
\item
  Let \(\mathcal{G}^{\mathrm{T}} = \mathcal{G}^{\mathrm{T}}(\mathfrak{p}(\mathrm{B}))\) be the inertia group of \(\mathfrak{p}(\mathrm{B})\) (Chapter V, § 2, no. 2) and let \(\mathbf{K}^{\mathrm{T}}\) denote the subfield of \(\mathbf{K}'\) consisting of the invariants of \(\mathcal{G}^{\mathrm{T}}\) and \(\mathbf{B}^{\mathrm{T}}\) the ring of the valuation induced on \(\mathbf{K}^{\mathrm{T}}\) by \(v\) , which is equal to \(\mathbf{B} \cap \mathbf{K}^{\mathrm{T}}\) ; recall that, if \(k\) and \(k'\) are the residue fields of A and B, the residue field \(k^{\mathrm{T}}\) of \(\mathbf{B}^{\mathrm{T}}\) is the greatest separable extension of \(k\) contained in the quasi-Galois extension \(k'\) of \(k\) and that \(\mathcal{G}^{\mathrm{Z}} / \mathcal{G}^{\mathrm{T}}\) is canonically isomorphic to the group of \(k\) -automorphisms of \(k'\) (or \(k^{\mathrm{T}}\) ) (Chapter V, § 2, Theorem 2 and Proposition 5). Let \(v^{\mathrm{T}}\) be the restriction of \(v\) to \(\mathbf{K}^{\mathrm{T}}\) ; show that the order group of \(v^{\mathrm{T}}\) is also equal to \(\Gamma\) (apply Lemma 2 of no. 1 to the extension \(\mathbf{K}^{\mathrm{T}}\) of \(\mathbf{K}^{\mathrm{Z}}\) ).
\end{enumerate}

\begin{enumerate}
\def\labelenumi{\arabic{enumi}.}
\setcounter{enumi}{8}
\tightlist
\item
  Let K be a field, L a finite algebraic extension of K, \(v'\) a valuation on L, A' its ring, v the restriction of \(v'\) to K and \(A = A' \cap K\) its ring.
\end{enumerate}

\begin{enumerate}
\def\labelenumi{(\alph{enumi})}
\item
  Suppose that A is Henselian (Exercise 6); show that the product \(e(v'/v)f(v'/v)\) divides \(n = [\mathbf{L}:\mathbf{K}]\) and that the quotient \(n / e(v'/v)f(v'/v)\) is a power of the characteristic exponent of the residue field \(k\) of \(v\) . (Reduce it to the case where L is a Galois extension of K; use Exercises 6 (e) and 7 and Chapter V, §2, Theorem 2 and Proposition 5.)
\item
  Suppose further that \(n\) is not divisible by the characteristic exponent of the residue field \(k\) of \(v\) and that \(f(v' / v) = 1\) . Show that, for every integer \(m\) equal to the order of an element \(\gamma\) of the group \(\Gamma_{v'} / \Gamma_v\) there exists \(x \in \mathbf{L}\) such that the class of \(v'(x) \mod \Gamma_v\) is equal to \(\gamma\) and \(x^m \in \mathbf{K}\) .
\end{enumerate}

¶ 10. Let w be the 2-adic valuation on the 2-adic field \(Q_{2}\) ; on the field of rational functions \(\mathbf{Q}_2(\mathbf{X})\) consider the discrete valuation \(v\) extending \(w\) and such that

\[
v (a _ {0} + a _ {1} \mathrm{X} + \dots + a _ {n} \mathrm{X} ^ {n}) = \inf _ {i} (w (a _ {i}))
\]

(§ 10, no. 1, Lemma 1); let K be the completion of \(\mathbf{Q}_2(\mathbf{X})\) with respect to this valuation and let its valuation also be denoted by \(v\) . Let L be the splitting field of the polynomial \((\mathrm{Y}^2 - \mathrm{X})^2 - 2\) in K{[}Y{]} and let \(v'\) be the unique valuation on L extending \(v\) . Show that {[}L: K{]} = 8, \(e(v'/v) = 4\) , \(f(v'/v) = 2\) ; if \(k\) (resp. \(k'\) ) is the residue field of \(v\) (resp. \(v'\) ), \(k'\) is a radical extension of \(k\) of degree 2; show that there exists no sub-extension E of L such that {[}E: K{]} = 2 for which the residue field of the restriction of \(v'\) to E is equal to \(k'\) . (Of necessity \(\mathrm{E} = \mathrm{K}(\sqrt{\alpha})\) , where \(\alpha \in \mathrm{K}\) would not be a square in K but such that \(\alpha \equiv \mathrm{X}\) (mod. 2); express \(\sqrt{\alpha}\) with the aid of a suitable basis of L over \(\mathrm{K}(\sqrt{2})\) and observe that there exists in \(\mathrm{K}(\sqrt{2})\) no element whose square is congruent to X mod. 2.)

¶ 11. Let K be a field, L a finite Galois extension of K and G its Galois group. Let v be a valuation on L with residue field k and order group Γ; suppose that v is invariant under G and that the restriction v\textbar K has the same residue field k as v, so that \(G = G^{Z} = G^{T}\) (notation of Exercise 8).

\begin{enumerate}
\def\labelenumi{(\alph{enumi})}
\item
  Let \(x \in \mathbf{L}^*\) , \(\sigma \in \mathcal{G}\) . Show that \(x^{-1}\sigma(x)\) is a unit for \(v\) , that its image \(\varepsilon_{\sigma}(x)\) in \(k^*\) depends only on the valuation \(v(x)\) and the class of \(\sigma\) modulo the commutator group \(\mathcal{G}'\) of \(\mathcal{G}\) and that there is thus defined a Z-bilinear mapping (also denoted by \(\varepsilon\) ) \((\mathcal{G}/\mathcal{G}') \times \Gamma \to k^*\) , which is equal to 1 on \((\mathcal{G}/\mathcal{G}') \times v(\mathbf{K}^*)\) .
\item
  For all \(\sigma \in \mathcal{G}\) , let \(\phi(\sigma)\) be the homomorphism \(\Gamma / v(\mathbf{K}^*) \to k^*\) mapping \(\varepsilon_{\sigma}(x)\) to \(v(x)\) for all \(x \in \mathbf{L}^*\) ; thus there is defined a homomorphism
\end{enumerate}

\[
\phi \colon \mathcal {G} \to \operatorname{Hom} _ {\mathbf {Z}} (\Gamma / v (\mathrm{K} ^ {*}), k ^ {*}).
\]

Show that, if \(p\) is the characteristic exponent of \(k\) , the kernel \(\mathcal{N}\) of \(\phi\) has order a power of \(p\) . (In the contrary case, there would exist a \(\sigma \neq e\) in \(\mathcal{N}\) and a prime number \(q \neq p\) such that \(\sigma^q = e\) ; for some integer \(x\) in \(L - K\) , write \(\sigma(x) = x + y\) , where \(v(y) > v(x)\) , and calculate \(\sigma^q(x)\) to obtain a contradiction.) Deduce that \(\mathcal{G}\) is a solvable group.

\begin{enumerate}
\def\labelenumi{(\alph{enumi})}
\setcounter{enumi}{2}
\tightlist
\item
  Suppose that \(n = [\mathbf{L} : \mathbf{K}]\) is prime to \(p\) ; show that \(\phi\) is bijective, that \(e(\mathbf{L} / \mathbf{K}) = n\) and that, if \(\nu\) is the \(lcm\) of the orders of the elements of \(\mathcal{G}\) , \(k^*\) contains the \(\nu\) -th roots of unity. (Use (b) and Lemma 2 of no. 1.)
\end{enumerate}

¶ 12. Let K be a field, v a valuation on K, A the ring of the valuation v and Γ its order group.

\begin{enumerate}
\def\labelenumi{(\alph{enumi})}
\tightlist
\item
  Let \(f(\mathbf{X}) = a_{0}\mathbf{X}^{n} + a_{1}\mathbf{X}^{n-1} + \cdots + a_{n}\) be a polynomial in A{[}X{]} such that \(v(a_{0}) = 0\) , all of whose roots belong to K; let \(x_{i} (1 \leqslant i \leqslant r)\) be these distinct roots and let \(k_{i}\) be the order of multiplicity of \(x_{i} (1 \leqslant i \leqslant r)\) . Let
\end{enumerate}

\[
g (\mathbf {X}) = b _ {0} \mathbf {X} ^ {n} + b _ {1} \mathbf {X} ^ {n - 1} + \dots + b _ {n} = b _ {0} (\mathbf {X} - y _ {1}) \dots (\mathbf {X} - y _ {n})
\]

another polynomial in A{[}X{]} such that \(v(b_{0}) = 0\) and the \(y_{h}\) belong to K. Show that there exists \(\lambda_{0} \in \Gamma\) such that \(v(x_{i} - x_{j}) < \lambda_{0}\) for every ordered pair of distinct indices i, j and, for all \(\lambda \geqslant \lambda_{0}\) , there exists \(\mu \geqslant \lambda\) such that the relations \(v(a_{i} - b_{i}) \geqslant \mu\) for all i imply that, for \(1 \leqslant i \leqslant r\) , there are exactly \(k_{i}\) indices h such that \(v(x_{i} - y_{i}) \geqslant \lambda\) . (First evaluate \(v(f(y_{h}))\) in two ways to show that necessarily \(v(x_{i} - y_{h}) \geqslant \lambda\) for at least one i; then evaluate

\[
v \left(\frac {f ^ {\prime} (z)}{f (z)} - \frac {g ^ {\prime} (z)}{g (z)}\right)
\]

at suitable points \(z \in \mathbf{K}\) .) If further \(b_{i} = a_{i}\) except for \(i = n - 1\) , show that, so long as \(\lambda_0\) is sufficiently large, the \(y_{h}\) are all distinct (evaluate \(f'(y_h)/f(y_h)\) assuming that \(y_{h}\) is not a simple root of \(g\) ).

\begin{enumerate}
\def\labelenumi{(\alph{enumi})}
\setcounter{enumi}{1}
\tightlist
\item
  Suppose henceforth that \(\mathbf{K}\) is the algebraic closure of a subfield \(\mathbf{K}_0\) such that \(\mathbf{A}_0 = \mathbf{A} \cap \mathbf{K}_0\) is Henselian. Suppose that the polynomial \(f\) belongs to \(\mathbf{A}_0[\mathbf{X}]\) and is irreducible and separable over \(\mathbf{K}_0\) . Let \(z \in \mathbf{K}\) be such that
\end{enumerate}

\[
v (z - v _ {i}) > v (z - x _ {j})
\]

for all \(j \neq i\) ; then show that \(\mathrm{K}_0(x_i) \subset \mathrm{K}_0(z)\) . (Show that \(z - x_i\) is equal to all its conjugates with respect to \(\mathrm{K}_0(x_i)\) .)

\begin{enumerate}
\def\labelenumi{(\alph{enumi})}
\setcounter{enumi}{2}
\tightlist
\item
  Suppose that \(f \in \mathbf{A}_0[\mathbf{X}]\) is separable and let \(f = a_0 \prod_{i=1}^{\infty} f_i\) be the decomposition of \(f\) into monic irreducible polynomials in \(\mathbf{K}_0[\mathbf{X}]\) (which belong in fact to \(\mathbf{A}_0[\mathbf{X}]\) , cf.~Chapter V, § 1, no. 3, Proposition 11). Show (in the notation of (a)) that, if \(\mu\) is taken sufficiently large, the decomposition of \(g\) into monic
\end{enumerate}

irreducible factors in \(K_{0}[X]\) can be written as \(b_{0}\prod_{i=1}^{n}g_{i}\) , where, for all i, \(g_{i}\) is of the same degree as \(f_{i}\) (a decomposition said to be ``of the same type'' as that of f) and further the splitting fields of \(f_{i}\) and \(g_{i}\) are the same for all i. (Show first that g is separable; consider then a finite Galois extension of \(K_{0}\) containing the roots of f and g and consider the way in which the Galois group of this extension permutes these roots; finally use (b).)

\begin{enumerate}
\def\labelenumi{(\alph{enumi})}
\setcounter{enumi}{3}
\tightlist
\item
  Give an example where \(f\) is irreducible but not separable and where, for all \(\mu \geqslant \lambda_0\) , there exists a polynomial
\end{enumerate}

\[
g (\mathbf {X}) = b _ {0} \mathbf {X} ^ {n} + b _ {1} \mathbf {X} ^ {n - 1} + \dots + b _ {n} \in \mathrm{A} _ {0} [ \mathbf {X} ]
\]

where \(v(a_{i} - b_{i}) \geqslant \mu\) for all i, which is not irreducible (take \(K_{0}\) such that \(K \neq K_{0}\) and K is a radial extension of \(K_{0}\) ).

¶ 13. (a) Let K be a field and v a non-improper valuation on K. Let E be a set filtered by an ultrafilter U and \(\xi \mapsto x_{\xi}\) , \(\xi \mapsto y_{\xi}\) two mappings of E to K with the topology \(\mathcal{T}_{v}\) . Show that, if the mapping \(\xi \mapsto x_{\xi}y_{\xi}\) converges to 0 in K with respect to U, so does one of the mappings \(\xi \mapsto x_{\xi}\) , \(\xi \mapsto y_{\xi}\) . (Observe that, if \(\xi \mapsto x_{\xi}\) has no cluster point with respect to \(\mathfrak{U}\) , the mapping \(z \mapsto x_{\xi}^{-1}\) is bounded on a subset of \(\mathfrak{U}\) .)

\begin{enumerate}
\def\labelenumi{(\alph{enumi})}
\setcounter{enumi}{1}
\item
  Let \(f\) be a non-constant polynomial in \(\mathbf{K}[\mathbf{X}]\) . Show that, if \(\xi \mapsto f(x_{\xi})\) converges to 0 in \(\mathbf{K}\) with respect to \(\mathfrak{U}\) , \(\xi \mapsto x_{\xi}\) converges in \(\hat{\mathbf{K}}\) (decompose \(f\) into factors in an algebraic extension of \(\mathbf{K}\) and use (a)).
\item
  Suppose that K is algebraically closed in \(\hat{\mathbf{K}}\) . Then the mapping \(x \mapsto f(x)\) of K to itself is closed (use (b)). Deduce that, if \(g\) is a rational function in \(\mathbf{K}(\mathbf{X})\) and B is a bounded closed subset of K, then \(g(\mathbf{B} - \mathbf{P})\) (where P is the set of poles of \(g\) in B) is closed in K.
\item
  Suppose that the algebraic closure of \(\mathbf{K}\) is a radicial extension of \(\mathbf{K}\) . Then show that \(\hat{\mathbf{K}}\) is algebraically closed. (Apply (b) to an algebraic closure of \(\hat{\mathbf{K}}\) and a polynomial \(f \in \hat{\mathbf{K}}[\mathbf{X}]\) ; use also Exercise 12 (a).)
\end{enumerate}

¶ 14. (a) For a field K to be Henselian with a valuation v, it is necessary and sufficient that K be the greatest separable algebraic extension of K contained in \(\hat{K}\) and that \(\hat{K}\) be Henselian. (To see that the condition is necessary, use Exercise 12 (b); to show that it is sufficient, observe that, to verify condition (H) of Chapter III, § 4, Exercise 3, we may confine our attention to the case where P is separable over K, noting that, if Q is an irreducible factor of P in \(\hat{K}[X]\) such that \(Q^{p^{e}}\) belongs to K{[}X{]}, then Q is a g.c.d. of P and \(Q^{p^{e}}\) in \(\hat{K}[X]\) .) Consider the case where v is of height 1 (cf.~Exercise 6 (b)).

\begin{enumerate}
\def\labelenumi{(\alph{enumi})}
\setcounter{enumi}{1}
\item
  Deduce from (a) that there exists in the \(p\) -adic field \(\mathbf{Q}_p\) a countable subfield which is Henselian with the \(p\) -adic valuation.
\item
  Give an example of a field K which is Henselian with a discrete valuation which is not complete and such that \(\hat{K}\) is a finite radical extension of K (cf.~Chapter V, § 1, Exercise 20).
\end{enumerate}

¶ 15. (a) Let K be a field, \(v_1\) , \(v_2\) two independent valuations (§ 7, no. 2) on K and \(K_1\) , \(K_2\) the completions of K with respect to \(v_1\) and \(v_2\) respectively. Let \(L_1\) , \(L_2\) be two fields such that \(K \subset L_1 \subset K_1\) , \(K \subset L_2 \subset K_2\) , which are Henselian (with the extensions by continuity of \(v_1\) and \(v_2\) respectively). Let \(g_1 \in L_1[X]\) , \(g_2 \in L_2[X]\) be two separable polynomials with the same degree n.~Show that there exists a polynomial \(h \in K[X]\) which, in \(L_i[X]\) , has the same type of decomposition (Exercise 12 (c)) as \(g_i\) ( \(i = 1, 2\) ). (Reduce it to the case where \(g_1\) and \(g_2\) are in K{[}X{]} and are monic; consider the polynomial

\[
a ^ {n} g _ {1} (\mathbf {X} / a) + b ^ {n} g _ {2} (\mathbf {X} / a) - \mathbf {X} ^ {n},
\]

where \(v_{1}(a) \leqslant 0\) , \(v_{2}(a) > 0\) is arbitrarily large and \(v_{1}(b) > 0\) is arbitrarily large.)

\begin{enumerate}
\def\labelenumi{(\alph{enumi})}
\setcounter{enumi}{1}
\item
  Deduce from (a) that, if K is Henselian with \(v_{1}\) , the algebraic closure of \(L_{2}\) is radial over \(L_{2}\) and \(K_{2}\) is algebraically closed (take \(g_{2}\) to be irreducible and \(g_{1} \in K[X]\) to be the product of n distinct factors of prime degree).
\item
  Deduce from (b) that, if \(\mathbf{K}\) is Henselian with \(v_{1}\) and \(v_{2}\) , the algebraic closure of \(\mathbf{K}\) is radicial over \(\mathbf{K}\) .
\item
  Show that, if K is Henselian with a discrete valuation v, the algebraic closure of K cannot be radicial over K and therefore K cannot be Henselian with any valuation independent of v (use Algebra, Chapter V, § 11, Exercise 12).
\end{enumerate}

\begin{enumerate}
\def\labelenumi{\arabic{enumi}.}
\setcounter{enumi}{15}
\tightlist
\item
  Let K be a Henselian field with a valuation v of height 1. If L is an infinite separable algebraic extension of K, show that L cannot be complete with the valuation extending v (and also denoted by v). (Form a sequence \((x_{p})\) of elements of L such that the degree \(n_{p}\) of \(x_{p}\) with respect to K tends to \(+\infty\) and \(v(x_{p+1}-x_{p})\) is strictly greater than the valuations of the differences between \(x_{p}\) and its conjugates with respect to K; show that the sequence \((x_{p})\) is a Cauchy sequence which cannot converge in E, using Exercise 12 (b).)
\end{enumerate}

¶ 17. Let K be a Henselian field with a valuation v which is not algebraically closed, and L a subfield of K such that {[}K: L{]} is finite. Show that under these conditions L is Henselian with the restriction of v. (In the contrary case, there would exist two extensions \(v_{1}, v_{2}\) of \(v|L\) to an algebraic closure \(\Omega\) of K such that the restriction of \(v_{1}\) and \(v_{2}\) to a finite algebraic extension \(L'\) of L would not be equivalent. Deduce that there would exist on a suitable finite quasi-Galois extension \(K'\) of L containing K, two valuations \(v_{1}', v_{2}'\) which are not equivalent and for which \(K'\) would be Henselian; using Exercise 15 (c), show that, if \(E = L^{p^{-\infty}}\) (p being the characteristic exponent of \(\Omega\) ), \(E \neq \Omega\) but \(\Omega\) would be a finite extension of E; complete the argument with the aid of Algebra, Chapter VI, § 2, Exercise 31.)

¶ 18. Let K be a Henselian field with a valuation v, k the residue field of v and suppose that every finite algebraic extension of k is cyclic over k (which is for example the case when k is finite). Let A be the ring of v and \(f \in A[X]\) a monic polynomial such that, if \(f(\mathbf{X}) = (\mathbf{X} - \alpha_{1}) \ldots (\mathbf{X} - \alpha_{n})\) , where the \(\alpha_{i}\) belong to an algebraic closure of K, the element \(D = \prod_{i < j} (\alpha_{i} - \alpha_{j})^{2}\) of A (``discriminant'' of f) is such that \(v(D) = 0\) . Let \(\bar{f}_{j} (1 \leqslant j \leqslant s)\) be the irreducible factors of the canonical image \(\bar{f}\) of f in k{[}X{]} and let \(r_{j}\) be the degree of \(\bar{f}_{j}\) . Show that the Galois over K of the splitting field of f, considered as a group of permutations of the \(\alpha_{i}\) is generated by a permutation \(\sigma\) which decomposes into cycles of respective lengths \(r_{1}, r_{2}, \ldots, r_{s}\) (observe that there exists only a single extension of k of given degree, up to a k-isomorphism). Deduce that, for D to be a square in A, it is necessary and sufficient that n - s be even (``Stickelberger's Theorem''; examine under what condition D is invariant under \(\sigma\) ).

\begin{enumerate}
\def\labelenumi{\arabic{enumi}.}
\setcounter{enumi}{18}
\tightlist
\item
  Let K be a Henselian field with a valuation v. Let E be a finite-dimensional vector space over K and Q a non-degenerate quadratic form on E such that the relation \(Q(x) = 0\) implies x = 0. Let
\end{enumerate}

\[
\Phi (x, y) = \mathrm{Q} (x + y) - \mathrm{Q} (x) - \mathrm{Q} (y)
\]

be the associated bilinear form. Show that \(2v(\Phi(x,y)) \geqslant v(Q(x)) + v(Q(y))\) ; deduce that

\[
v (\mathrm{Q} (x + y)) \geqslant \inf (v (\mathrm{Q} (x)), v (\mathrm{Q} (y))).
\]

¶ 20. Let K be a field of characteristic ≠2 which is Henselian with a valuation v; let ξ ↦ ξ̄ be an involutive automorphism of K such that v(ξ̄) = v(ξ) and let Φ be a non-degenerate Hermitian form on a finite-dimensional vector space E over K.

\begin{enumerate}
\def\labelenumi{(\alph{enumi})}
\item
  Let \(U = (\alpha_{ij})\) be the matrix of \(\Phi\) with respect to a basis \((e_i)\) of E; show that there exists an element \(\lambda\) of the order group of \(v\) such that, if \(\Phi'\) is another Hermitian form on E whose matrix \(U' = (\alpha_{ij}')\) with respect to \((e_i)\) satisfies the conditions \(v(\alpha_{ij} - \alpha_{ij}') > \lambda\) for every ordered pair \((i,j)\) , then \(\Phi\) and \(\Phi'\) are equivalent. (Argue as in Algebra, Chapter IX, § 6, Exercise 6, using Exercise 14 (a) above.)
\item
  Suppose that \(\xi \mapsto \overline{\xi}\) is the identity (hence \(\Phi\) is a symmetric form) and that the valuation \(v\) is discrete, normed and such that \(v(2) = 0\) ; let \(\pi\) be a uniformizer for \(v\) . Show that \(\Phi\) is characterized, up to equivalence, by its index \(\nu\) and by two symmetric bilinear forms \(\Psi_1, \Psi_2\) of index 0 on vector spaces \(k^r\) and \(k^s\) respectively, where \(k\) is the residue field of \(v\) and \(r + s = n - 2\) . (With the aid of a Witt decomposition, reduce it to the case where \(\Phi\) is of index 0; with the aid of Exercise 19, show that, for all \(i \geqslant 0\) , the set \(\mathbf{M}_i\) of \(x \in \mathbf{E}\) such that \(v(\Phi(x, x)) \geqslant i\) is a module over the ring A of the valuation \(v\) . Show that, if \(x \in \mathbf{M}_i\) , \(y \in \mathbf{M}_{i+1}\) , then \(v(\Phi(x, y)) \geqslant i + 1\) and, by taking quotients, symmetric bilinear forms can therefore be derived from \(\Phi\) on the vector \(k\) -spaces \(\mathbf{M}_0 / \mathbf{M}_1\) and \(\mathbf{M}_1 / \mathbf{M}_2\) . Use Exercise 6 (c) of §3 and the fact that the equation \(\xi^2 = \alpha\) has a solution in A for \(\alpha \equiv 1 (\text{mod. } \pi)\) .) Consider the case where \(k\) is finite (cf.~Algebra, Chapter IX, §6, Exercise 4).
\end{enumerate}

\begin{enumerate}
\def\labelenumi{\arabic{enumi}.}
\setcounter{enumi}{20}
\tightlist
\item
  Let K be a field, v a discrete valuation on K, A the ring of v, \(\pi\) a uniformizer of v and k the residue field of v.
\end{enumerate}

\begin{enumerate}
\def\labelenumi{(\alph{enumi})}
\item
  Let P, R be two polynomials in A{[}X{]}, P being monic, and suppose that: (1) \(\deg(R) < h \cdot \deg(P)\) , where \(h \geqslant 1\) is an integer; (2) the canonical image \(\overline{P}\) of P in k{[}X{]} is irreducible. Show then that, if the polynomial \(Q = P^{h} + \pi R\) is reducible in \(\widehat{K}[X]\) , necessarily h \textgreater{} 1 and \(\overline{P}\) divides the canonical image \(\overline{R}\) of R in k{[}X{]}. Deduce that, for a given monic irreducible polynomial \(r(X) \in k[X]\) and a given integer \(h \geqslant 1\) , there exists a separable irreducible polynomial \(Q \in A[X]\) such that its canonical image \(\overline{Q}\) is equal to \(r^{h}\) .
\item
  Let \(k' = k(\alpha)\) be an algebraic extension of \(k\) of degree \(m\) and \(h\) an integer \(\geqslant 1\) . Show that there exists an algebraic extension \(L\) of \(K\) of degree \(hm\) such that there is only a single valuation \(v'\) (up to equivalence) on \(L\) extending \(v\) , that \(e(v'/v) = h\) , \(f(v'/v) = m\) and that the residue field of \(v'\) is isomorphic to \(k'\) (use (a)).
\end{enumerate}

\begin{enumerate}
\def\labelenumi{\arabic{enumi}.}
\setcounter{enumi}{21}
\tightlist
\item
  \begin{enumerate}
  \def\labelenumii{(\alph{enumii})}
  \tightlist
  \item
    If there exists a discrete valuation on a field \(\mathbf{K}\) , show that the algebraic closure of \(\mathbf{K}\) is infinite over \(\mathbf{K}\) .
  \end{enumerate}
\end{enumerate}

\begin{enumerate}
\def\labelenumi{(\alph{enumi})}
\setcounter{enumi}{1}
\tightlist
\item
  Let K be a finitely generated extension of a field \(K_{0}\) . Show that, if K is not algebraic over \(K_{0}\) , there exists a discrete valuation v on K such that \(v(x) = 0\) on \(K_{0}\) .
\end{enumerate}

\exercisesection{9}

\begin{enumerate}
\def\labelenumi{\arabic{enumi}.}
\tightlist
\item
  \begin{enumerate}
  \def\labelenumii{(\alph{enumii})}
  \tightlist
  \item
    Let K be a field (not necessarily commutative) and T a locally compact topology on K compatible with the ring structure on K; show that T is compatible with the field structure on K. (Use the Theorem of R. Ellis (General Topology, Chapter X, § 3, Exercise 25) or argue directly reproducing the proofs of Integration, Chapter VII, § 1, no. 10 and those of Proposition 1 of this paragraph.)
  \end{enumerate}
\end{enumerate}

\begin{enumerate}
\def\labelenumi{(\alph{enumi})}
\setcounter{enumi}{1}
\tightlist
\item
  Give an example of a locally compact topology on a commutative field K, which is compatible with the additive group structure on K, but not with its ring structure. (Take K to be the field of fractions of a compact integral domain A (A being for example a ring of formal power series \(k[[X, Y]]\) , where k is finite) and take as fundamental system of neighbourhoods of 0 in K the neighbourhoods of 0 in A.)
\end{enumerate}

¶ 2. (a) In a space \(\mathbf{R}^n\) ( \(n \geqslant 2\) ), let U be a non-empty open set with non-empty complement; show that the frontier of U contains a non-empty perfect set (cf.~General Topology, Chapter I, § 9, Exercise 17).

\begin{enumerate}
\def\labelenumi{(\alph{enumi})}
\setcounter{enumi}{1}
\item
  Deduce from (a) that, if \(\mathbf{A}\) is an everywhere dense subset of \(\mathbf{R}^n\) which meets every perfect set in \(\mathbf{R}^n\) , \(\mathbf{A}\) is convex.
\item
  Show that there exists in \(\mathbf{C}\) an everywhere dense subfield \(\mathbf{K}\) which is connected and locally connected and is a pure transcendental extension of \(\mathbf{Q}\) (apply General Topology, Chapter IX, § 5, Exercise 18 (b) and (c), constructing \(\mathbf{K}\) by transfinite induction, using (b) and the method described in Set Theory, Chapter III, § 6, Exercise 24). Deduce that there exists a subfield \(\mathbf{K}' \supset \mathbf{K}\) of \(\mathbf{C}\) which is (algebraically) isomorphic to \(\mathbf{R}\) , connected and locally connected.
\item
  Show, using (c), that there exists on C the topology of a connected and locally connected space, which is compatible with the field structure and for which the completion of C is an algebra over C which is the direct composition of two fields isomorphic to C.
\end{enumerate}

\begin{enumerate}
\def\labelenumi{\arabic{enumi}.}
\setcounter{enumi}{2}
\tightlist
\item
  Let K be a totally disconnected non-discrete locally compact commutative field; let A be the ring of the absolute value mod \(_{K}\) on K and let U be the group of units of A. For every integer n \textgreater{} 0, let \(_{n}U\) denote the group of n-th roots of unity in K and \(U^{n}\) the subgroup of U consisting of the n-th powers of the elements of U. Show that, if n is prime to the characteristic of K, \(U^{n}\) is an open subgroup of U and that
\end{enumerate}

\[
\operatorname{Card} (\mathrm{U} / \mathrm{U} ^ {n}) = \operatorname{mod} _ {\mathrm{K}} (n). \operatorname{Card} (_ {n} \mathrm{U})
\]

(use Exercise 14 of Integration, Chapter VII, §2 and Commutative Algebra, Chapter III, §4, no. 6, Corollary 1 to Theorem 2 to show that, if m is the maximal ideal of A, the image under \(x \mapsto x^n\) of \(1 + \mathfrak{m}^k\) is \(1 + n \cdot \mathfrak{m}^k\) for \(k\) sufficiently large).

\begin{enumerate}
\def\labelenumi{\arabic{enumi}.}
\setcounter{enumi}{3}
\tightlist
\item
  \begin{enumerate}
  \def\labelenumii{(\alph{enumii})}
  \tightlist
  \item
    Let K be a commutative field and v a valuation on K such that K is Henselian with v (§ 8, Exercise 6). Suppose further that K and the residue field k of v are of characteristic 0. Show that there exists a subfield \(K_{0}\) of the ring A of v such that the canonical mapping \(A \rightarrow k\) , restricted to \(K_{0}\) , is an isomorphism of \(K_{0}\) onto k. (Let H be a subfield of A such that the image of H under the canonical mapping is an isomorphism onto a subfield E of k; show that, if \(E \neq k\) , there exists \(\alpha \notin H\) in A such that the subfield \(\mathrm{H}(\alpha)\) of K is contained in A and is canonically isomorphic to \(\mathrm{E}(\bar{\alpha})\) , where \(\bar{\alpha}\) is the class of \(\alpha\) in k; distinguish two cases according to whether \(\bar{\alpha}\) is algebraic or transcendental over E.)
  \end{enumerate}
\end{enumerate}

\begin{enumerate}
\def\labelenumi{(\alph{enumi})}
\setcounter{enumi}{1}
\tightlist
\item
  Suppose further that v is a discrete valuation and that K is complete with respect to v; deduce from (a) that K is isomorphic to the field of formal power series \(k((\mathbf{T}))\) .
\end{enumerate}

\begin{enumerate}
\def\labelenumi{\arabic{enumi}.}
\setcounter{enumi}{4}
\tightlist
\item
  \begin{enumerate}
  \def\labelenumii{(\alph{enumii})}
  \tightlist
  \item
    Suppose that K is a commutative field, v a valuation of height 1 on K such that K is complete with respect to v and A the ring of v; suppose further that the residue field k of v is perfect and of characteristic p \textgreater{} 0. For every element \(\xi \in k^{*}\) and every integer n, let \(x_{n}\) be an element of the class \(\xi^{p-n}\) in A; show that the sequence \((x_{n}^{p^{n}})\) is a Cauchy sequence in A whose limit is independent of the choice of the \(x_{n}\) in the classes \(\xi^{p-n}\) . If this limit is denoted by \(\phi(\xi)\) , show that \(\phi\) is the unique isomorphism u of the multiplicative group \(k^{*}\) to the multiplicative group \(K^{*}\) such that, for all \(\xi \in k^{*}\) , \(u(\xi)\) is an element of A belonging to the class \(\xi\) .
  \end{enumerate}
\end{enumerate}

\begin{enumerate}
\def\labelenumi{(\alph{enumi})}
\setcounter{enumi}{1}
\item
  If \(\mathbf{K}\) is also of characteristic \(p\) , show that \(\phi\) , extended to \(k\) by \(\phi(0) = 0\) , is an isomorphism of the field \(k\) onto a subfield of \(\mathbf{K}\) . Deduce a new proof of Theorem 1 (iii) of no. 3.
\item
  Suppose that \(k\) is finite. Show that, if \(r\) is prime to \(p\) , the group \((\mathbf{K}^*)^r\) of \(r\) -th powers of elements of \(\mathbf{K}^*\) is of finite index in \(\mathbf{K}^*\) (use Hensel's Lemma). Show that, if further \(v\) is discrete and \(\mathbf{K}\) of characteristic 0, the same result is valid without restriction on \(r\) (observe that every element of \(1 + p^2\mathbf{A}\) is a \(p\) -th power).
\end{enumerate}

\exercisesection{10}

\begin{enumerate}
\def\labelenumi{\arabic{enumi}.}
\tightlist
\item
  Let K be a field and P the prime subfield of K; the absolute dimension of K is the number dim. al \(_{P}\) K if P is of characteristic p \textgreater{} 0 and the number dim. al \(_{P}\) K + 1 if P is of characteristic 0. Let v be a valuation on K, h its height, r its rational rank and k its residue field.
\end{enumerate}

\begin{enumerate}
\def\labelenumi{(\alph{enumi})}
\item
  Suppose that the absolute dimension \(n\) of \(K\) is finite. Then, if \(s\) is the absolute dimension \(k\) , \(r + s \leqslant n\) .
\item
  Suppose further that K is a finitely generated extension of P. Then, if \(r + s = n\) , k is a finitely generated extension of its prime field and the order group of v is isomorphic to \(Z'\) ; if \(h + s = n\) , k is a finitely generated extension of its prime field and the order group of v is isomorphic to \(Z'\) ordered lexicographically; finally, if s = n - 1, v is a discrete valuation and k is a finitely generated extension of its prime field.
\end{enumerate}

¶ 2. Let K be a field, v a valuation on K, L an extension of K and \(v'\) a valuation on L extending v; L is called an immediate extension (with respect to \(v'\) ) if \(e(v'/v) = f(v'/v) = 1\) . The completion \(\hat{K}\) of K is an immediate extension of K.

\begin{enumerate}
\def\labelenumi{(\alph{enumi})}
\item
  For \(\mathbf{L}\) to be an immediate extension of \(\mathbf{K}\) , it is necessary and sufficient that, for all \(x \in \mathbf{L} - \mathbf{K}\) , there exist \(y \in \mathbf{K}\) such that \(v'(x - y) > v'(x)\) .
\item
  For every field K, show that there exists a maximal immediate extension L of K, that is, not admitting an immediate extension distinct from itself (use § 3, Exercise 3 (b)).
\item
  Show that, if K is linearly compact with the topology defined by v, it admits no immediate extension other than itself (use (a), noting that, in the notation of (a), the set of \(v'(x - y)\) for \(y \in K\) has no greatest element in the value group of \(v'\) ).
\item
  Suppose that \(\mathbf{K}\) is not linearly compact; let \(\mathbf{B}\) be a well-ordered subset of the order group of \(v\) and \((a_{\beta})_{\beta \in \mathbb{B}}\) a family of elements of \(\mathbf{K}\) such that, for \(\lambda < \mu < \nu\) ,
\end{enumerate}

\[
v (a _ {\lambda} - a _ {\mu}) <   v (a _ {\mu} - a _ {v}),
\]

but that there exists no \(x \in \mathbf{K}\) such that \(v(x - a_{\lambda}) = v(a_{\lambda} - a_{\mu})\) for every ordered pair \((\lambda, \mu)\) such that \(\lambda < \mu\) . For every extension \(\mathbf{E}\) of \(\mathbf{K}\) and every extension \(w\) of \(v\) to \(\mathbf{E}\) , let \(\mathrm{D}_{\mathbf{E}}(\lambda)\) denote the set of \(z \in \mathbf{E}\) such that \(v(z - a_{\lambda}) \geqslant \gamma_{\lambda}\) , where \(\gamma_{\lambda}\) is the common value of the \(v(a_{\lambda} - a_{\mu})\) for \(\lambda < \mu\) ; the hypothesis is therefore that \(\bigcap_{\beta \in \mathbf{B}} \mathrm{D}_{\mathbf{K}}(\beta) = \varnothing\) (§ 5, Exercise 5 (a)). Let \(\Omega\) be an algebraic closure of \(\mathbf{K}\) and let \(v_0\) be an extension of \(v\) to \(\Omega\) ; show that, if \(\mathbf{P}\) is a polynomial in \(\mathbf{K}[\mathbf{X}]\) , for there to exist \(\lambda \in \mathbf{B}\) such that \(v(\mathrm{P}(a_{\mu})) = v(\mathrm{P}(a_{\lambda}))\) for all \(\lambda \leqslant \mu\) , it is necessary and sufficient that no zero of \(\mathbf{P}\) in \(\Omega\) belong to \(\bigcap_{\beta \in \mathbf{B}} \mathrm{D}_{\Omega}(\beta)\) . If \(\mathbf{P}\) has this property and \(\mathbf{E}\) and \(w\) have the same meaning as above, show that, for all \(z \in \bigcap_{\beta \in \mathbf{B}} \mathrm{D}_{\mathbf{E}}(\beta), w(\mathrm{P}(z) - \mathrm{P}(a_{\mu})) > w(\mathrm{P}(z)) = v(\mathrm{P}(a_{\mu}))\) so long as \(\mu\) is large enough (decompose \(\mathrm{P}(\mathbf{X})\) into factors in \(\Omega[\mathbf{X}]\) ). Deduce that one of the following situations holds:

\begin{enumerate}
\def\labelenumi{(\arabic{enumi})}
\tightlist
\item
  Either \(\bigcap_{\beta \in B} D_{\Omega}(\beta) \neq \varnothing\) and, if \(\theta\) is an element of this intersection whose degree over \(K\) is the least possible, \(K(\theta)\) is an immediate extension of \(K\) distinct from \(K\) .
\end{enumerate}

\begin{enumerate}
\def\labelenumi{(\arabic{enumi})}
\setcounter{enumi}{1}
\tightlist
\item
  Or \(\bigcap_{\beta \in B} D_{\Omega}(\beta) = \varnothing\) ; then there is a valuation \(v'\) on K(X) extending \(v\) such that \(v(P(X)) = v(P(a_\mu))\) for \(\mu\) large enough and K(X), with \(v'\) , is an immediate extension of K (use the criterion in (a)).
\end{enumerate}

\begin{enumerate}
\def\labelenumi{(\alph{enumi})}
\setcounter{enumi}{4}
\tightlist
\item
  Deduce from (c) and (d) that for K to be linearly compact, it is necessary and sufficient that K have no immediate extension other than itself.
\end{enumerate}

\chapter{Divisors}

All the rings considered in this chapter are assumed to be commutative and to possess a unit element. All the ring homomorphisms are assumed to map unit element to unit element. Every subring of a ring A is assumed to contain the unit element of A.

\section{KRULL DOMAINS}

\subsection{Divisorial ideals of an integral domain}

DEFINITION 1. Let A be an integral domain and K its field of fractions. Every sub-A-module a of K such that there exists an element \(d \neq 0\) in A for which da ⊂ A is called a fractional ideal of A (or of K, by an abuse of language).

Every finitely generated sub-A-module \(\mathfrak{a}\) of \(\mathbf{K}\) is a fractional ideal: for if \((a_{i})_{1\leqslant i\leqslant n}\) is a system of generators of \(\mathfrak{a}\) , we may write \(a_{i} = b_{i} / d_{i}\) , where \(b_{i}\in \mathbf{A}\) , \(d_{i}\in \mathbf{A}\) and \(d_{i}\neq 0\) ; if \(d = d_{1}\ldots d_{n}\) , clearly \(da\subset A\) . In particular the monogenous sub-A-modules of \(\mathbf{K}\) are fractional ideals (recall that they have been called fractional principal ideals in Algebra, Chapter VI, § 1, no. 5). If \(\mathbf{A}\) is Noetherian, every fractional ideal is a finitely generated A-module. Every sub-A-module of a fractional ideal of \(\mathbf{A}\) is a fractional ideal. Every ideal of \(\mathbf{A}\) is a fractional ideal; to avoid confusion, these will also be called the integral ideals of \(\mathbf{A}\) .

We denote by I(A) the set of non-zero fractional ideals of A. Given two elements a, b of I(A), we shall write \(a \prec b\) (or \(b \succ a\) ) for the relation ``every fractional principal ideal containing a also contains b''; clearly this relation is a preordering on I(A). Let R denote the associated equivalence relation ``a \(\prec b\) and \(b \prec a\) '' (Set Theory, Chapter III, § 1, no. 2) and D(A) the quotient set I(A)/R; we shall say that the elements of D(A) are the divisors of A and, for every fractional ideal \(a \in I(A)\) , we shall denote by div a (or div \(_{A}\) a) the canonical image of a in D(A) and we shall say that div a is the divisor of a; if a = Ax is a fractional principal ideal, we write \(\operatorname{div}(x)\) instead of \(\operatorname{div}(\mathbf{A}x)\) and \(\operatorname{div}(x)\) is called the divisor of \(x\) ; the elements of D(A) of the form \(\operatorname{div}(x)\) are called principal divisors. By taking the quotient, the preordering \(\prec\) on I(A) defines on D(A) an ordering which we shall denote by \(\leqslant\) .

For all \(\mathfrak{a} \in \mathrm{I}(\mathbf{A})\) there exists by hypothesis some \(d \neq 0\) in \(\mathbf{A}\) such that \(\mathfrak{a} \subset \mathrm{Ad}^{-1}\) ; the intersection \(\tilde{\mathfrak{a}}\) of the fractional principal ideals containing \(\mathfrak{a}\) is therefore an element of \(\mathrm{I}(\mathbf{A})\) . Clearly the relation \(\mathfrak{a} \prec \mathfrak{b}\) is equivalent to the relation \(\tilde{\mathfrak{a}} \supset \tilde{\mathfrak{b}}\) ; the relation \(\mathfrak{a} \supset \mathfrak{b}\) therefore implies \(\mathfrak{a} \prec \mathfrak{b}\) . For two elements \(\mathfrak{a}, \mathfrak{b}\) of \(\mathrm{I}(\mathbf{A})\) to be equivalent modulo \(\mathbb{R}\) , it is necessary and sufficient that \(\tilde{\mathfrak{a}} = \tilde{\mathfrak{b}}\) .

DEFINITION 2. Every element \(\mathfrak{a}\) of I(A) such that \(\mathfrak{a} = \tilde{\mathfrak{a}}\) is called a divisorial fractional ideal of A.

In other words a divisorial ideal is just a non-zero intersection of a non-empty family of fractional principal ideals. Every non-zero intersection of divisorial ideals is a divisorial ideal. If a is divisorial, so is ax for all \(x \in K^{*}\) , the mapping \(b \mapsto bx\) being a bijection of the set of fractional principal ideals onto itself. For all \(a \in I(A)\) , \(\tilde{a}\) is the least divisorial ideal containing a and is equivalent to a modulo R. Moreover, if b is a divisorial ideal equivalent to a modulo R, then \(\tilde{a} = \tilde{b} = b\) . Hence \(\tilde{a}\) is the unique divisorial ideal b such that div a = div b (in other words, the restriction of the mapping \(a \mapsto div a\) to the set of divisorial ideals is injective).

Let \(\mathfrak{a}\) and \(\mathfrak{b}\) be two fractional ideals of \(\mathbf{K}\) . Recall (Chapter I, § 2, no. 10) that \(\mathfrak{b}\) : \(\mathfrak{a}\) denotes the set of \(x \in \mathbf{K}\) such that \(x\mathfrak{a} \subset \mathfrak{b}\) ; this is obviously an A-module; if \(\mathfrak{b} \in \mathrm{I}(\mathrm{A})\) and \(\mathfrak{a} \in \mathrm{I}(\mathrm{A})\) , then \(\mathfrak{b} \colon \mathfrak{a} \in \mathrm{I}(\mathrm{A})\) ; for if \(d\) is a non-zero element of \(\mathbf{A}\) such that \(db \subset \mathbf{A}\) and \(da \subset \mathbf{A}\) and \(a\) is a non-zero element of \(\mathbf{A} \cap \mathfrak{a}\) , then \(da(\mathfrak{b} : \mathfrak{a}) \subset \mathbf{A}\) ; on the other hand, if \(b \neq 0\) belongs to \(\mathfrak{b}\) , then \(bda \subset \mathfrak{b}\) , hence \(bd \in \mathfrak{b}\) : \(\mathfrak{a}\) and \(\mathfrak{b} : \mathfrak{a} \neq 0\) .

The definition of b: a can also be written:

\[
\mathfrak {b}: \mathfrak {a} = \bigcap_ {x \in \mathfrak {a}, x \neq 0} \mathfrak {b} x ^ {- 1}.\tag{1}
\]

PROPOSITION 1. (a) If \(\mathfrak{b}\) is a divisorial ideal and \(\mathfrak{a} \in \mathbf{I}(\mathbf{A})\) , \(\mathfrak{b}\) : \(\mathfrak{a}\) is divisorial.

\begin{enumerate}
\def\labelenumi{(\alph{enumi})}
\setcounter{enumi}{1}
\item
  Let \(\mathfrak{a},\mathfrak{b}\) be in I(A). In order that \(\operatorname {div}\mathfrak{a} = \operatorname {div}\mathfrak{b}\) , it is necessary and sufficient that A: \(\mathfrak{a} = \mathbf{A}\) : \(\mathfrak{b}\) .
\item
  For all \(\mathfrak{a} \in \mathbf{I}(\mathbf{A})\) , \(\tilde{\mathfrak{a}} = \mathbf{A} : (\mathbf{A} : \mathfrak{a})\) .
\end{enumerate}

Assertion (a) follows immediately from equation (1) since, if \(\mathfrak{b}\) is divisorial, so is \(\mathfrak{bx}^{-1}\) for all \(x \neq 0\) .

To show (b), let \(\mathrm{P}(\mathfrak{a})\) denote the set of fractional principal ideals containing \(\mathfrak{a}\) ; the relation \(Ax \in \mathrm{P}(\mathfrak{a})\) is equivalent to \(x^{-1}\mathfrak{a} \subset A\) and hence to \(x^{-1} \in A\) : \(\mathfrak{a}\) . As the relation \(\operatorname{div} \mathfrak{a} = \operatorname{div} \mathfrak{b}\) is by definition equivalent to \(\mathrm{P}(\mathfrak{a}) = \mathrm{P}(\mathfrak{b})\) , it is also equivalent to \(A\) : \(\mathfrak{a} = A\) : \(\mathfrak{b}\) .

Finally, as \(\mathfrak{a}(\mathbf{A}:\mathfrak{a})\subset \mathbf{A},\mathfrak{a}\subset \mathbf{A}:(\mathbf{A}:\mathfrak{a})\) . Replacing \(\mathfrak{a}\) by \(\mathbf{A}:\mathfrak{a}\) in this formula, it is seen that \(\mathbf{A}:\mathfrak{a}\subset \mathbf{A}:(\mathbf{A}:(\mathbf{A}:\mathfrak{a}))\) ; on the other hand, the relation \(\mathfrak{a}\subset \mathbf{A}:(\mathbf{A}:\mathfrak{a})\) implies

\[
\mathrm{A}: \mathfrak {a} \supset \mathrm{A}: (\mathrm{A}: (\mathrm{A}: \mathfrak {a})).
\]

Therefore A: a = A : (A : (A : a)) and it follows from (b) that div a = div(A : (A : a)) As A : (A : a) is divisorial by (a), certainly \(\tilde{a} = A\) : (A : a), which proves (c).

Remark. In the course of the above proof it has been proved that A: a = A: (A: (A: a)) for every ideal \(a \in I(A)\) , which is a special case of Set Theory, Chapter III, § 1, no. 5, Proposition 2.

PROPOSITION 2. (i) In D(A) every non-empty set bounded above admits a least upper bound. More precisely, if \((\mathfrak{a}_t)\) is a non-empty family of elements of I(A) which is bounded above, then

\[
\sup _ {i} (\operatorname{div} a _ {i}) = \operatorname{div} \left(\bigcap_ {i} \tilde {a} _ {i}\right).
\]

\begin{enumerate}
\def\labelenumi{(\roman{enumi})}
\setcounter{enumi}{1}
\tightlist
\item
  In D(A) every non-empty set bounded below admits a greatest lower bound. More precisely, if \((\mathfrak{a}_{\mathfrak{t}})\) is a non-empty family of elements of I(A) which is bounded below, then
\end{enumerate}

\[
\inf _ {i} (\operatorname{div} a _ {i}) = \operatorname{div} \left(\sum_ {i} a _ {i}\right).
\]

\begin{enumerate}
\def\labelenumi{(\roman{enumi})}
\setcounter{enumi}{2}
\tightlist
\item
  The set D(A) is a lattice.
\end{enumerate}

Let \((\mathfrak{a}_i)\) be a non-empty family of elements of I(A) which is bounded above. To say that a divisorial ideal \(\mathfrak{b}\) bounds this family above amounts to saying that it is contained in all the \(\tilde{\mathfrak{a}}_i\) , that is that \(\mathfrak{b}\) is contained in \(\bigcap_{i} \tilde{\mathfrak{a}}_i\) . Hence \(\bigcap_{i} \tilde{\mathfrak{a}}_i \neq (0)\) and \(\bigcap_{i} \tilde{\mathfrak{a}}_i\) is therefore a divisorial ideal, which shows (i).

Now let \((\tilde{a}_t)\) be a non-empty family of elements of I(A) which is bounded below. To say that a divisorial ideal \(b\) bounds this family below means that it contains all the \(\tilde{a}_t\) , that is (since \(b\) is divisorial) that it contains all the \(a_t\) , or also that \(b \supset \sum_{i} a_i\) . This proves (ii).

Finally, to prove (iii) it is sufficient by (i) and (ii) to prove that, if a, b are in I(A), the set \(\{a, b\}\) is bounded both above and below in I(A); now it is bounded above by \(a \cap b\) (which is distinct from (0)). It is bounded below by \(a + b\) , for \(a + b \in I(A)\) : if d and \(d'\) are non-zero elements of A such that \(da \subset A\) and \(d'b \subset A\) , then \(dd'(\mathfrak{a} + \mathfrak{b}) \subset \mathbf{A}\) .

COROLLARY. If \(x, y\) and \(x + y\) are in \(\mathbf{K}^*\) , then \(\operatorname{div}(x + y) \geqslant \inf(\operatorname{div}(x), \operatorname{div}(y))\) . \(\mathbf{A}(x + y) \subset \mathbf{A}x + \mathbf{A}y\) and hence \(\operatorname{div}(x + y) \geqslant \operatorname{div}(\mathbf{A}x + \mathbf{A}y)\) .
